\section{Statistik und Statistical Learning}

% ============================================================
\subsection{Skalenniveaus und Charakterisierung von Merkmalen}
% ============================================================

\textbf{Datentypen:}
\textbf{Strukturierte Daten} haben eine vordefinierte Form (Tabellen mit Zeilen und Spalten);
\textbf{unstrukturierte Daten} besitzen keine formalisierte Struktur (Texte, Bilder, Audio, Sensordaten).

\medskip
\noindent Die \textbf{Datenmatrix} ist eine $n \times p$-Matrix, deren Zeilen die \emph{Merkmalsträger} (Beobachtungen, $i=1,\ldots,n$) und deren Spalten die \emph{Merkmale} (Variablen, $j=1,\ldots,p$) darstellen.

\medskip
\noindent\textbf{Klassifikation von Variablen:}

\begin{center}
\small
\begin{tabular}{|l|l|}
\hline
\textbf{Typ} & \textbf{Beschreibung / Beispiel} \\
\hline
Qualitativ (kategorial) & Endliche Kategorien, z.\,B. Geschlecht, Augenfarbe \\
Quantitativ (numerisch) & Zahlenwerte, z.\,B. Alter, Größe \\
Diskret & Endlich/abzählbar viele Ausprägungen, z.\,B. Personenanzahl \\
Stetig & Unendlich viele Werte im Intervall, z.\,B. Körpergewicht \\
Manifest & Direkt beobachtbar, z.\,B. Körpergröße \\
Latent & Nur über Indikatoren messbar, z.\,B. Arbeitszufriedenheit \\
Unabhängig & Einflussgröße (Prädiktor) \\
Abhängig & Zielgröße (Kriterium) \\
\hline
\end{tabular}
\end{center}

\medskip
\noindent\textbf{Skalenniveaus} bilden eine Hierarchie, die bestimmt, welche mathematischen Operationen und Lagemaße zulässig sind:

\begin{center}
\small
\begin{tabular}{|l|l|l|l|}
\hline
\textbf{Skala} & \textbf{Eigenschaft} & \textbf{Operationen} & \textbf{Lagemaße} \\
\hline
Nominal & reine Klassifikation & $=$, $\neq$ & Modus \\
Ordinal & Rangordnung, kein fixer Abstand & $=$, $\neq$, $<$, $>$ & Modus, Median \\
Intervall & feste Abstände, kein wahrer Nullpunkt & $+$, $-$ zusätzlich & Modus, Median, $\bar{x}$ \\
Verhältnis & wahrer Nullpunkt & $\cdot$, $\div$ zusätzlich & Modus, Median, $\bar{x}$ \\
Absolut & natürliche Einheit (Stückzahl) & alle & Modus, Median, $\bar{x}$ \\
\hline
\end{tabular}
\end{center}

\noindent Beispiele: \emph{Nominal} -- Nationalität; \emph{Ordinal} -- Schulnoten, Likert-Skalen; \emph{Intervall} -- Temperatur in °C, IQ; \emph{Verhältnis} -- Länge, Preis; \emph{Absolut} -- Anzahl von Personen.

% ============================================================
\newpage

\subsection{Lage-, Streuungs- und Assoziationsmaße}
% ============================================================

\textbf{Ziel:} Erste Eindrücke über die Verteilung von Daten durch kompakte Kennzahlen gewinnen.

\medskip
\noindent\textbf{Grafische Darstellung -- Univariat:}

\begin{center}
\small
\begin{tabular}{|l|l|}
\hline
\textbf{Diagrammtyp} & \textbf{Einsatz} \\
\hline
Balkendiagramm & Kategoriale und metrische Häufigkeiten \\
Kreisdiagramm & Kategoriale Daten mit wenigen Kategorien (Anteile) \\
Histogramm & Kontinuierliche Verteilungen (Klassen, Dichte) \\
Boxplot & Median, Quartile, Ausreißer; ordinal und metrisch \\
Violin-Plot & Boxplot + Kernel-Density-Schätzung (Verteilungsform) \\
\hline
\end{tabular}
\end{center}

\medskip
\noindent\textbf{Grafische Darstellung -- Bivariat:}

\begin{center}
\small
\begin{tabular}{|l|l|}
\hline
\textbf{Diagrammtyp} & \textbf{Einsatz} \\
\hline
Streudiagramm (Scatterplot) & 2 metrische Variablen, Korrelation/Trends \\
Boxplot nach Gruppen & Metrisch vs. kategorial, Vergleich von Lagen \\
Gruppiertes Balkendiagramm & 2 kategoriale oder gemischte Variablen \\
Liniendiagramm & Zeitreihen, Trends über Kategorien \\
Heatmap & Korrelationsmatrizen, 2-D-Intensitäten \\
\hline
\end{tabular}
\end{center}

\noindent\textbf{Grundsatz:} Maximale Informationsdichte bei minimaler visueller Komplexität.
Die Wahl des Diagrammtyps hängt von der Art der Variablen (kategorial/metrisch), der Anzahl der Variablen und dem Analyseziel (Häufigkeiten, Verteilungsform, Zusammenhang) ab.

% ============================================================
\newpage

\subsection{Ähnlichkeitsmaße und Distanzmaße}
% ============================================================

\textbf{Grundkonzept:} \emph{Distanz} $d$ misst Unähnlichkeit, \emph{Ähnlichkeit} $s$ ist deren Inverses.
Bei normierten Maßen gilt: $s = 1 - d$.

\medskip
\noindent\textbf{Distanzmaße für numerische Daten:}

\begin{itemize}[leftmargin=1.5em]
    \item \textbf{Euklidische Distanz:}
    $d_{\text{Euk}} = \sqrt{\sum_{j=1}^{p}(x_{B,j}-x_{A,j})^2}$ --
    geometrischer Luftlinienabstand; empfindlich gegenüber Ausreißern und unterschiedlichen Skalierungen (Standardisierung empfohlen).

    \item \textbf{Manhattan-Distanz:}
    $d_{\text{Man}} = \sum_{j=1}^{p}|x_{B,j}-x_{A,j}|$ --
    Summe der absoluten Koordinatendifferenzen; robuster gegenüber Ausreißern.

    \item \textbf{Minkowski-Distanz:}
    $d_k = \bigl(\sum_{j=1}^{p}|x_{B,j}-x_{A,j}|^k\bigr)^{1/k}$ --
    Verallgemeinerung mit Parameter $k$: $k=1$ Manhattan, $k=2$ Euklidisch, $k\to\infty$ Chebyshev ($\max_j|x_{A,j}-x_{B,j}|$).

    \item \textbf{Mahalanobis-Distanz:}
    $d_{\text{Mah}} = \sqrt{(\mathbf{x}_B-\mathbf{x}_A)^T\hat{\boldsymbol{\Sigma}}^{-1}(\mathbf{x}_B-\mathbf{x}_A)}$ --
    berücksichtigt Korrelationsstruktur und Skalierungen; erfordert invertierbare Kovarianzmatrix (instabil bei $p>n$).
\end{itemize}

\noindent\textbf{Ähnlichkeitsmaße:}

\begin{itemize}[leftmargin=1.5em]
    \item \textbf{Kosinus-Ähnlichkeit:}
    $\cos(\theta)=\frac{\mathbf{x}_A \cdot \mathbf{x}_B}{\|\mathbf{x}_A\|_2\|\mathbf{x}_B\|_2}\in[-1,1]$ --
    misst den Winkel zwischen Vektoren unabhängig von ihrer Länge; Standard für Text- und Bildvektoren (TF-IDF, Word Embeddings).

    \item \textbf{Tanimoto-Koeffizient (Jaccard):}
    $S = \frac{|A \cap B|}{|A \cup B|}\in[0,1]$ -- für Binärvektoren/Mengen; Anwendung in Genetik und Chemoinformatik.

    \item \textbf{Dice-Koeffizient:}
    $S_{\text{Dice}}=\frac{2|A\cap B|}{|A|+|B|}$ -- gewichtet die Überschneidung stärker als Tanimoto; Zusammenhang: $S_{\text{Dice}}=\frac{2S_{\text{Tan}}}{1+S_{\text{Tan}}}$; Einsatz in NLP und Biologie.

    \item \textbf{Pearson-/Spearman-Distanz:}
    $d = 1-r$ bzw. $d = 1-r_S$ -- korrelationsbasiert; Spearman ist robust gegenüber Ausreißern und monoton nichtlinearen Zusammenhängen.
\end{itemize}

\noindent\textbf{Faustregeln zur Wahl:}
Euklidisch -- normalisierte numerische Daten $\cdot$
Mahalanobis -- korrelierte Variablen $\cdot$
Kosinus -- hochdimensionale Text-/Bildvektoren $\cdot$
Tanimoto/Dice -- Binärvektoren $\cdot$
Spearman -- Ausreißer/Rangdaten $\cdot$
Manhattan -- allgemeine Robustheit.

% ============================================================
\newpage

\subsection{Diagrammtypen und Einsatzgebiete}
% ============================================================

\noindent\textit{Hinweis: Dieses Core Topic ist in den vorliegenden Ausarbeitungen nicht eigenständig ausgearbeitet, sondern wird in den Subsections „Lage-, Streuungs- und Assoziationsmaße" sowie „Ähnlichkeitsmaße und Distanzmaße" abgedeckt. Die folgende Zusammenfassung konsolidiert die relevanten Inhalte.}

\medskip
\noindent Die Wahl eines Diagrammtyps hängt von drei Faktoren ab: \emph{Skalenniveau} der Variablen, \emph{Anzahl} der betrachteten Variablen und dem \emph{Analyseziel} (Häufigkeiten, Verteilungsform, Zusammenhang, Zeitverlauf).

\medskip
\noindent\textbf{Übersicht nach Analyseziel:}

\begin{center}
\small
\begin{tabular}{|l|l|l|}
\hline
\textbf{Ziel} & \textbf{Variablentyp} & \textbf{Diagramm} \\
\hline
Häufigkeitsverteilung & Kategorial (nominal, ordinal) & Balken-, Kreisdiagramm \\
Häufigkeitsverteilung & Metrisch, stetig & Histogramm \\
Lage \& Streuung & Metrisch / ordinal & Boxplot \\
Verteilungsform & Metrisch & Violin-Plot, Histogramm + Dichte \\
Linearer Zusammenhang & 2 metrische Variablen & Streudiagramm (Scatterplot) \\
Gruppenvergleich & Metrisch vs. kategorial & Boxplot nach Gruppen \\
Kategorialer Vergleich & 2 kategoriale Variablen & Gestapeltes/gruppiertes Balkendiagramm \\
Zeitlicher Verlauf & Ordinal / metrisch & Liniendiagramm \\
Intensitätsmatrix & 2 Variablen (z.\,B. Korrelation) & Heatmap \\
\hline
\end{tabular}
\end{center}

\medskip
\noindent\textbf{Wichtige Eigenschaften einzelner Typen:}

\begin{itemize}[leftmargin=1.5em]
    \item \textbf{Boxplot:} Zeigt Median (Q2), unteres/oberes Quartil (Q1, Q3), Whiskers und Ausreißer. Kompakt für Vergleiche über Gruppen.
    \item \textbf{Violin-Plot:} Erweiterung des Boxplots mit Kernel-Density-Schätzung -- sichtbar, ob die Verteilung multimodal ist.
    \item \textbf{Scatterplot:} Ein Punkt pro Beobachtung $(x_i, y_i)$; Erweiterung als Bubble Chart mit dritter Variable als Punktgröße.
    \item \textbf{Heatmap:} Farben kodieren Wertebereiche einer zweidimensionalen Matrix, z.\,B. einer Korrelationsmatrix.
    \item \textbf{Kreisdiagramm:} Nur sinnvoll bei wenigen Kategorien; Menschen vergleichen Winkel schlechter als Längen (daher Balkendiagramm oft bevorzugt).
\end{itemize}

\noindent\textbf{Grundsatz:} Maximale Informationsdichte bei minimaler visueller Komplexität -- unnötige 3-D-Effekte und überladene Legenden vermeiden.

% ============================================================
\newpage

\subsection{Wahrscheinlichkeiten und Wahrscheinlichkeitsräume}
% ============================================================

\textbf{Zufallsexperiment:} Ein Prozess, dessen Ausgang nicht im Voraus bekannt ist, dessen mögliche Ergebnisse (Elementarereignisse $\omega$) aber vollständig bekannt sind. Genau ein $\omega$ tritt mit Sicherheit ein.

\medskip
\noindent\textbf{Wahrscheinlichkeitsinterpretationen:}
\begin{itemize}[leftmargin=1.5em]
    \item \textbf{Frequentistisch:} Relative Häufigkeit eines Ereignisses bei sehr häufiger Wiederholung unter identischen Bedingungen -- objektiv, beobachterunabhängig.
    \item \textbf{Bayessch:} Subjektiver Grad der Gewissheit basierend auf verfügbarem Vorwissen; ermöglicht Aussagen wie „99\,\% Wahrscheinlichkeit, dass der Fingerabdruck von Person X stammt".
\end{itemize}

\medskip
\noindent\textbf{Formale Grundstruktur -- Wahrscheinlichkeitsraum $(\Omega, \mathcal{A}, P)$:}

\begin{itemize}[leftmargin=1.5em]
    \item $\Omega$ \textbf{(Stichprobenraum/Ergebnisraum):} Menge aller möglichen Ausgänge; kann endlich, abzählbar oder überabzählbar sein.
    \item $\mathcal{A}$ \textbf{($\sigma$-Algebra/Ereignisraum):} Menge aller Ereignisse $E \subseteq \Omega$, denen Wahrscheinlichkeiten zugeordnet werden. Eigenschaften: (i) $\Omega \in \mathcal{A}$; (ii) $E\in\mathcal{A} \Rightarrow E^c\in\mathcal{A}$; (iii) abzählbare Vereinigungen von Ereignissen $\in\mathcal{A}$.
    \item $P$ \textbf{(Wahrscheinlichkeitsfunktion):} $P:\mathcal{A}\to[0,1]$ nach den \textbf{Kolmogorov-Axiomen}:
    \begin{enumerate}[leftmargin=1.5em]
        \item $P(\Omega)=1$
        \item $P(E)\geq 0$ für alle $E\in\mathcal{A}$
        \item $\sigma$-Additivität: $P\!\left(\bigcup_{i=1}^\infty E_i\right)=\sum_{i=1}^\infty P(E_i)$ für paarweise disjunkte $E_i$
    \end{enumerate}
\end{itemize}

\noindent\textbf{Laplace'scher Wahrscheinlichkeitsraum:} Endlicher Raum mit gleichverteilten Elementarereignissen (Indifferenzprinzip). Für $m$ mögliche Fälle und $g$ günstige Fälle gilt:
\[
    P(E) = \frac{g}{m} = \frac{\text{Anzahl günstiger Fälle}}{\text{Anzahl möglicher Fälle}}
\]
\textit{Beispiel Würfel:} $\Omega=\{1,\ldots,6\}$, Ereignis „gerade Zahl" $E=\{2,4,6\}$: $P(E)=\frac{3}{6}=\frac{1}{2}$.

\medskip
\noindent\textbf{Abgeleitete Rechenregeln:}
\begin{itemize}[leftmargin=1.5em]
    \item \textbf{Komplementärregel:} $P(E^c) = 1 - P(E)$
    \item \textbf{Additionsregel:} $P(E_1 \cup E_2) = P(E_1) + P(E_2) - P(E_1 \cap E_2)$
    \item \textbf{Monotonie:} $E_1 \subseteq E_2 \Rightarrow P(E_1) \leq P(E_2)$
\end{itemize}





% ============================================================

\subsection{Kombinatorik}
% ============================================================

\definition{\textbf{Kombinatorik:} Zählen von Anordnungen und Auswahlen endlicher Mengen. Fundamental für Laplace-Räume: $P(E) = \frac{\text{günstige Fälle}}{\text{mögliche Fälle}}$.}

\begin{center}
\small
\begin{tabular}{|l|c|c|}
\hline
\textbf{Typ} & \textbf{Ohne Wiederholung} & \textbf{Mit Wiederholung} \\
\hline
\textbf{Permutation} (alle $n$ Elemente anordnen) & $n!$ & $\dfrac{n!}{n_1!\cdots n_k!}$ \\[6pt]
\hline
\textbf{Variation} ($k$ aus $n$, Reihenfolge zählt) & $\dfrac{n!}{(n-k)!}$ & $n^k$ \\[6pt]
\hline
\textbf{Kombination} ($k$ aus $n$, Reihenfolge zählt nicht) & $\dbinom{n}{k} = \dfrac{n!}{k!(n-k)!}$ & $\dbinom{n+k-1}{k}$ \\[6pt]
\hline
\end{tabular}
\end{center}

\textbf{Entscheidungsregel:}
\begin{itemize}[leftmargin=1.5em]
    \item Werden \emph{alle} $n$ Objekte angeordnet? $\to$ \textbf{Permutation}
    \item Werden nur $k < n$ Objekte ausgewählt? Reihenfolge wichtig? $\to$ \textbf{Variation}; nicht wichtig? $\to$ \textbf{Kombination}
    \item Wiederholungen erlaubt? $\to$ rechte Spalte
\end{itemize}

\textbf{Eigenschaften des Binomialkoeffizienten:} $\binom{n}{k} = \binom{n}{n-k}$ (Symmetrie); $\binom{n}{0} = \binom{n}{n} = 1$; $\binom{n}{k} + \binom{n}{k+1} = \binom{n+1}{k+1}$ (Pascal'sche Identität).

\textbf{Beispiele:}
Permutation: 3 Personen auf $3! = 6$ Arten anordnen.
Permutation mit WH: Buchstaben in „AABBCC": $\frac{6!}{2!\cdot2!\cdot2!} = 90$.
Variation ohne WH: Erste 3 Plätze aus 10 Pferden: $\frac{10!}{7!} = 720$.
Variation mit WH: 3-stellige Codes aus Ziffern 0--9: $10^3 = 1000$.
Kombination ohne WH: 5 Karten aus 52: $\binom{52}{5} = 2\,598\,960$.
Kombination mit WH: 3 Früchte aus 5 Sorten: $\binom{7}{3} = 35$.

% ============================================================
\newpage

\subsection{Diskrete und stetige Zufallsvariablen}
% ============================================================

\textit{[Dieses Core Topic war in beiden vorliegenden Ausarbeitungen nicht inhaltlich ausgearbeitet -- bitte Inhalt nachliefern.]}

% ============================================================
\newpage

\subsection{Statistische Tests}
% ============================================================

\definition{\textbf{Hypothesentest:} Verfahren zur Entscheidung, ob eine Vermutung über einen Parameter oder eine Verteilung der Grundgesamtheit anhand von Stichprobendaten haltbar ist.}

\textbf{Hypothesenarten:}
\begin{itemize}[leftmargin=1.5em]
    \item \textbf{$H_0$ (Nullhypothese):} Konservative Standardannahme, meist „kein Effekt".
    \item \textbf{$H_1$ (Alternativhypothese):} Die zu prüfende Gegenhypothese.
    \item \textbf{Einfach vs. zusammengesetzt:} Einzelner Parameterwert ($H_0: \mu = 10$) vs. Wertebereich ($H_0: \mu \leq 10$).
    \item \textbf{Einseitig vs. zweiseitig:} $H_1: \mu > \mu_0$ bzw. $H_1: \mu < \mu_0$ vs. $H_1: \mu \neq \mu_0$. Einseitige Tests nur bei klar vorherbestimmter Effektrichtung gerechtfertigt.
\end{itemize}

\textbf{Ablauf eines Hypothesentests:}
\begin{enumerate}[leftmargin=1.5em]
    \item $H_0$, $H_1$ und Signifikanzniveau $\alpha$ festlegen.
    \item Teststatistik $t_{\text{obs}}$ aus der Stichprobe berechnen (z.\,B. $z$-, $t$-, $\chi^2$-, $F$-Statistik).
    \item Kritischen Bereich (Ablehnungsbereich) basierend auf $\alpha$ und der Verteilung von $T$ unter $H_0$ bestimmen.
    \item Entscheidung: $t_{\text{obs}}$ im Ablehnungsbereich $\Rightarrow$ $H_0$ ablehnen; sonst $H_0$ nicht ablehnen.
\end{enumerate}

\definition{\textbf{p-Wert:} Wahrscheinlichkeit, unter $H_0$ einen Wert der Teststatistik zu beobachten, der mindestens so extrem ist wie $t_{\text{obs}}$ (in Richtung $H_1$):
$$p = P(T \geq t_{\text{obs}} \mid H_0 \text{ wahr})$$
Entscheidungsregel: $p < \alpha \Rightarrow H_0$ ablehnen (Ergebnis statistisch signifikant); $p \geq \alpha \Rightarrow H_0$ nicht ablehnen. \textbf{Achtung:} Der p-Wert ist \emph{nicht} die Wahrscheinlichkeit, dass $H_0$ wahr ist!}

\definition{\textbf{Fehlertypen und Teststärke:}}
\begin{center}
\small
\begin{tabular}{|l|c|c|}
\hline
 & $H_0$ tatsächlich wahr & $H_0$ tatsächlich falsch \\
\hline
$H_0$ nicht ablehnen & Korrekt ($1-\alpha$) & \textbf{Fehler II.\ Art} ($\beta$) \\
\hline
$H_0$ ablehnen & \textbf{Fehler I.\ Art} ($\alpha$) & Korrekt (Power $1-\beta$) \\
\hline
\end{tabular}
\end{center}

\begin{itemize}[leftmargin=1.5em]
    \item \textbf{Fehler I.\ Art} ($\alpha$): $H_0$ abgelehnt, obwohl wahr („Fehlalarm").
    \item \textbf{Fehler II.\ Art} ($\beta$): $H_0$ nicht abgelehnt, obwohl $H_1$ wahr („falsch negativ").
    \item \textbf{Teststärke (Power)} $= 1 - \beta$: Wahrscheinlichkeit, $H_0$ richtig abzulehnen.
    \item \textbf{Tradeoff:} $\alpha$ und $\beta$ sind negativ korreliert. Beide gleichzeitig reduzieren $\Rightarrow$ größere Stichprobe nötig.
\end{itemize}

% ============================================================
\newpage

\subsection{Population, Stichproben und Konfidenzintervalle}
% ============================================================

\definition{\textbf{Grundkonzept:} Die \textbf{Grundgesamtheit (Population)} ist die Menge aller Objekte mit unbekannter Verteilung. Die \textbf{Stichprobe} ist eine zufällig gezogene Teilmenge. \textbf{Inferenzprinzip:} Von der empirischen Stichprobenverteilung auf die wahre Populationsverteilung schließen.}

\textbf{Punktschätzer} schätzen einen Populationsparameter mit einem einzelnen Wert:

\begin{center}
\small
\begin{tabular}{|l|l|}
\hline
\textbf{Parameter} & \textbf{Schätzer} \\
\hline
Mittelwert $\mu$ & $\bar{X} = \frac{1}{n}\sum_{i=1}^n X_i$ \\[4pt]
\hline
Varianz $\sigma^2$ & $S^2 = \frac{1}{n-1}\sum_{i=1}^n (X_i - \bar{X})^2$ \\[4pt]
\hline
Anteil $p$ & $\hat{p} = \frac{\text{Anzahl Erfolge}}{n}$ \\[4pt]
\hline
\end{tabular}
\end{center}

\definition{\textbf{Konfidenzintervall (KI):} Ein zufälliges Intervall $[\theta_l, \theta_u]$ mit $P(\theta_l \leq \theta \leq \theta_u) = 1 - \alpha$. \textbf{Interpretation:} Bei wiederholter Stichprobenziehung enthält $(1-\alpha) \cdot 100\%$ aller so konstruierten Intervalle den wahren Parameter. Das Intervall ist zufällig; $\theta$ ist fest (aber unbekannt). Konfidenzniveaus: typisch $1-\alpha \in \{0{,}90;\, 0{,}95;\, 0{,}99\}$.}

Typen: \textbf{zweiseitig} $[\theta_l, \theta_u]$ mit $P(\theta<\theta_l) = P(\theta>\theta_u) = \frac{\alpha}{2}$; \textbf{einseitig unten} $[\theta_l,+\infty)$; \textbf{einseitig oben} $(-\infty, \theta_u]$.

\definition{\textbf{KI für $\mu$ bei bekannter Varianz $\sigma_0^2$:} Voraussetzung: Population normalverteilt oder $n \geq 30$ (CLT). Teststatistik: $Z = \frac{\bar{X} - \mu}{\sigma_0/\sqrt{n}} \sim N(0,1)$.
$$\text{KI: } \bar{X} \pm z_{1-\alpha/2} \cdot \frac{\sigma_0}{\sqrt{n}}$$
Für $\alpha = 0{,}05$: $z_{0{,}975} \approx 1{,}96$.}

\definition{\textbf{KI für $\mu$ bei unbekannter Varianz:} Schätzung von $\sigma^2$ durch $S^2$, Teststatistik: $T = \frac{\bar{X} - \mu}{S/\sqrt{n}} \sim t_{n-1}$.
$$\text{KI: } \bar{X} \pm t_{n-1,\, 1-\alpha/2} \cdot \frac{S}{\sqrt{n}}$$
Die $t$-Verteilung hat schwerere Schwänze als die Normalverteilung $\to$ breitere KIs, besonders bei kleinem $n$.}

\definition{\textbf{KI für $\sigma^2$ (normalverteilte Population):} Teststatistik: $Q = \frac{(n-1)S^2}{\sigma^2} \sim \chi^2_{n-1}$.
$$\text{KI: } \left[\frac{(n-1)S^2}{\chi^2_{n-1,\, 1-\alpha/2}},\quad \frac{(n-1)S^2}{\chi^2_{n-1,\, \alpha/2}}\right]$$}

% ============================================================
\newpage

\subsection{Zentraler Grenzwertsatz und Gesetz der großen Zahlen}
% ============================================================

\definition{\textbf{Unabhängig und identisch verteilt (u.i.v., iid):} Zufallsvariablen $X_1, \ldots, X_n$ heißen u.i.v., wenn sie unabhängig voneinander sind und alle dieselbe Verteilung wie $X$ haben (gleicher Erwartungswert $\mu$, gleiche Varianz $\sigma^2$). Die Zahlenwerte $x_1, \ldots, x_n$ heißen Realisierungen.}

\textbf{Hilfsungleichungen:}
\begin{itemize}[leftmargin=1.5em]
    \item \textbf{Markov-Ungleichung:} Für nicht-negative $X$ und $a > 0$: $P(X \geq a) \leq \frac{\mathbb{E}[X]}{a}$. Modellunabhängige obere Schranke.
    \item \textbf{Tschebyscheff-Ungleichung:} Für beliebiges $X$ mit Erwartungswert $\mu$ und Varianz $\sigma^2$:
    $$P(|X - \mu| \geq c) \leq \frac{\sigma^2}{c^2} \qquad \text{bzw.} \qquad P(|X - \mu| \geq k\sigma) \leq \frac{1}{k^2}$$
    Gilt \emph{unabhängig von der konkreten Verteilung}. Beispiele: $k=2 \to$ mind. 75\% im Intervall $[\mu-2\sigma, \mu+2\sigma]$; $k=3 \to$ mind. 88,9\%.
\end{itemize}

\definition{\textbf{Gesetz der großen Zahlen (LLN):} Für u.i.v. $X_1, \ldots, X_n$ mit endlichem $\mu$ und $\sigma^2$ gilt:
$$\bar{X}_n = \frac{1}{n}\sum_{i=1}^n X_i \xrightarrow{P} \mu \quad (n \to \infty)$$
d.\,h. $P(|\bar{X}_n - \mu| < c) \to 1$ für jedes $c > 0$. \textbf{Beweis via Tschebyscheff:} $\text{Var}(\bar{X}_n) = \frac{\sigma^2}{n} \to 0$, daher $P(|\bar{X}_n - \mu| \geq c) \leq \frac{\sigma^2}{nc^2} \to 0$. \textbf{Spezialfall (Bernoulli):} Die relative Häufigkeit $\frac{H_n}{n}$ konvergiert gegen die Wahrscheinlichkeit $p$.}

\definition{\textbf{Zentraler Grenzwertsatz (CLT):} Für u.i.v. $X_1, \ldots, X_n$ mit $\mathbb{E}[X] = \mu$ und $\text{Var}(X) = \sigma^2 > 0$ gilt:
$$Z_n = \frac{\bar{X}_n - \mu}{\sigma/\sqrt{n}} \xrightarrow{d} N(0,1) \quad \text{bzw.} \quad \bar{X}_n \xrightarrow{d} N\!\left(\mu,\, \frac{\sigma^2}{n}\right)$$
\textbf{Das Besondere:} Die Normalverteilung tritt \emph{unabhängig von der Verteilung} der $X_i$ auf -- gleichmäßig, exponentiell, diskret, etc. \textbf{Faustregel:} Wirkt ab $n \approx 30$, bei symmetrischen Verteilungen auch früher.}

\textbf{Praktische Relevanz:}
Viele reale Phänomene sind Summen vieler unabhängiger Einflussfaktoren (Messfehler, Körpergröße, Prüfungsnoten) und daher näherungsweise normalverteilt. Der CLT rechtfertigt die Konstruktion von Konfidenzintervallen und die Verwendung von $z$-Tests auch bei nicht-normalverteilten Populationen.

\textbf{Beispiel (Würfel):} $\mu = 3{,}5$, $\sigma^2 = 2{,}9$. Summe von $n=100$ Würfeln: $S_{100} \approx N(350, 290)$; Mittelwert: $\bar{X}_{100} \approx N(3{,}5;\, 0{,}029)$.

% ============================================================
\newpage

\subsection{Mehrdimensionale Zufallsvariablen}
% ============================================================

\definition{\textbf{Mehrdimensionaler Zufallsvektor:} $\mathbf{X} = (X_1, \ldots, X_n)^T$ mit Werten in $\mathbb{R}^n$. Im bivariaten Fall $(X, Y)$ interessiert die gemeinsame Struktur beider Variablen. Beispiel: $X$ = Körpergröße, $Y$ = Körpergewicht einer zufällig gewählten Person.}

\definition{\textbf{Gemeinsame Verteilung (Joint Distribution):}

Diskret: $f_{X,Y}(x,y) = P(X=x, Y=y)$, mit $\sum_x\sum_y f_{X,Y}(x,y) = 1$.

Stetig: Dichtefunktion $f_{X,Y}(x,y) \geq 0$ mit $\int\!\int f_{X,Y}(x,y)\,dx\,dy = 1$.

Gemeinsame CDF: $F_{X,Y}(x,y) = P(X \leq x, Y \leq y)$.

Die gemeinsame Verteilung enthält vollständige Information über beide Variablen und ihre Abhängigkeitsstruktur.}

\definition{\textbf{Randverteilungen:} Verteilung einer einzelnen Komponente durch Marginalisierung:
$$f_X(x) = \sum_y f_{X,Y}(x,y) \quad \text{(diskret)} \qquad f_X(x) = \int_{-\infty}^{\infty} f_{X,Y}(x,y)\,dy \quad \text{(stetig)}$$
Analog für $f_Y(y)$. Bei diskreten Variablen entsprechen die Randverteilungen den Zeilen-/Spaltensummen einer Kontingenztafel. Die Randverteilungen enthalten weniger Information als die gemeinsame Verteilung.}

\definition{\textbf{Bedingte Verteilungen:}
$$P(X=x \mid Y=y) = \frac{f_{X,Y}(x,y)}{f_Y(y)} \quad \text{(diskret)} \qquad f_{X|Y}(x\mid y) = \frac{f_{X,Y}(x,y)}{f_Y(y)} \quad \text{(stetig)}$$
Definiert für $f_Y(y) > 0$; die bedingte Dichte erfüllt $\int f_{X|Y}(x\mid y)\,dx = 1$.}

\definition{\textbf{Unabhängigkeit:} $X$ und $Y$ sind unabhängig, wenn:
$$f_{X,Y}(x,y) = f_X(x) \cdot f_Y(y) \quad \text{für alle } x, y$$
Äquivalent: Die bedingte Verteilung stimmt mit der Randverteilung überein. \textbf{Wichtig:} Unabhängigkeit $\Rightarrow$ Unkorreliertheit, aber \emph{nicht} umgekehrt.}

\begin{center}
\small
\begin{tabular}{|l|l|l|}
\hline
\textbf{Konzept} & \textbf{Diskret} & \textbf{Stetig} \\
\hline
Gemeinsam & $f_{X,Y}(x,y) = P(X=x, Y=y)$ & $f_{X,Y}(x,y)$ Dichtefunktion \\
\hline
Rand ($X$) & $f_X(x) = \sum_y f_{X,Y}(x,y)$ & $f_X(x) = \int f_{X,Y}(x,y)\,dy$ \\
\hline
Bedingt & $P(X=x|Y=y) = \frac{f_{X,Y}}{f_Y}$ & $f_{X|Y}(x|y) = \frac{f_{X,Y}}{f_Y}$ \\
\hline
Unabhängig & \multicolumn{2}{c|}{$f_{X,Y}(x,y) = f_X(x) \cdot f_Y(y)$} \\
\hline
\end{tabular}
\end{center}

% ============================================================
\newpage

\subsection{Kovarianz und Korrelation}
% ============================================================

\definition{\textbf{Kovarianz:} Maß für den linearen Zusammenhang zwischen $X$ und $Y$:
$$\text{Cov}(X,Y) = \sigma_{XY} = \mathbb{E}[(X-\mu_X)(Y-\mu_Y)] = \mathbb{E}[XY] - \mathbb{E}[X]\mathbb{E}[Y]$$
Stichprobe: $s_{XY} = \frac{1}{n-1}\sum_{i=1}^n (x_i - \bar{x})(y_i - \bar{y})$.
Vorzeichen: $>0$ positive, $<0$ negative Assoziation, $=0$ keine lineare Assoziation. \textbf{Nachteil:} Skalierungsabhängig -- Vergleiche zwischen Variablen mit unterschiedlichen Einheiten sinnlos.}

\definition{\textbf{Pearson-Korrelationskoeffizient:} Standardisiertes Maß, dimensionslos:
$$\rho_{X,Y} = \frac{\text{Cov}(X,Y)}{\sigma_X \cdot \sigma_Y} \in [-1,1] \qquad \text{Stichprobe: } r = \frac{s_{XY}}{s_X \cdot s_Y}$$
$\rho = \pm1$: perfekte lineare Beziehung; $\rho = 0$: keine lineare Beziehung (nichtlineare kann existieren). Faustregel: $|\rho| > 0{,}7$ stark; $0{,}4 < |\rho| \leq 0{,}7$ moderat; $|\rho| \leq 0{,}4$ schwach.}

\definition{\textbf{Eigenschaften:} Symmetrisch ($\rho_{X,Y} = \rho_{Y,X}$); invariant unter linearen Transformationen ($X' = aX+b$, $Y'= cY+d$ mit $a,c>0$); nur bei gemeinsamer Normalverteilung gilt $\rho = 0 \Leftrightarrow$ Unabhängigkeit.}

\textbf{Kovarianzmatrix und Korrelationsmatrix:} Für $\mathbf{X} = (X_1, \ldots, X_p)^T$:
$$\boldsymbol{\Sigma} = \mathbb{E}[(\mathbf{X}-\boldsymbol{\mu})(\mathbf{X}-\boldsymbol{\mu})^T], \quad \text{Diag: Varianzen, Off-Diag: Kovarianzen; symmetrisch, pos. semi-definit.}$$
Die \textbf{Korrelationsmatrix} $\mathbf{R}$ enthält $\rho_{ij}$ auf den Off-Diagonalen und 1 auf der Diagonalen. Visualisierung oft als Heatmap.

\definition{\textbf{Korrelation $\neq$ Kausalität:} Hohe Korrelation bedeutet nicht, dass $X$ die Ursache von $Y$ ist. Mögliche Szenarien:
\begin{itemize}[leftmargin=1.5em]
    \item Direkte Kausalität: $X \to Y$
    \item Umgekehrte Kausalität: $Y \to X$
    \item Gemeinsame Ursache (Confounding): Drittvariable $Z$ beeinflusst beide (Bsp.: Temperatur $\to$ Eiscreme-Konsum \emph{und} Ertrinken)
    \item Spurious Correlation: Zufälliger Zusammenhang, besonders bei großem $p$
\end{itemize}
\textbf{Simpson-Paradoxon:} Die Gesamtkorrelation kann ein anderes Vorzeichen haben als die Korrelation innerhalb von Subgruppen. Für Kausalitätsaussagen braucht man experimentelle Designs oder Causal Inference.}

% ============================================================
\newpage

\subsection{Regressionen}
% ============================================================

\definition{\textbf{Ziel:} Modellierung des Zusammenhangs zwischen Prädiktoren $x_1,\ldots,x_p$ und Response $y$ für Vorhersage, Inferenz über Effekte oder Modellierung $y = f(\mathbf{x}) + \varepsilon$.}

\subsubsection*{Einfache Lineare Regression (SLR)}

\definition{\textbf{Modell:} $y_i = \beta_0 + \beta_1 x_i + \varepsilon_i$ mit $\varepsilon_i \sim N(0, \sigma^2)$. \textbf{Annahmen:} Linearität, Unabhängigkeit der Fehler, Homoskedastizität (konstante Varianz), Normalverteilung der Fehler (für Inferenz).}

\textbf{OLS-Schätzer} (Minimierung von $\sum_i (y_i - \hat{y}_i)^2$):
$$\hat{\beta}_1 = \frac{s_{xy}}{s_x^2} = r_{xy}\frac{s_y}{s_x}, \qquad \hat{\beta}_0 = \bar{y} - \hat{\beta}_1\bar{x}$$

\textbf{Varianzzerlegung:} TSS $=$ ESS $+$ RSS, d.\,h. $\sum(y_i-\bar{y})^2 = \sum(\hat{y}_i-\bar{y})^2 + \sum e_i^2$.

\textbf{Bestimmtheitsmaß:} $R^2 = \frac{\text{ESS}}{\text{TSS}} = 1 - \frac{\text{RSS}}{\text{TSS}} \in [0,1]$. Anteil der erklärten Variabilität. Für SLR: $R^2 = r_{xy}^2$.

\subsubsection*{Multiple Lineare Regression (MLR)}

\definition{\textbf{Modell:} $\mathbf{y} = \mathbf{X}\boldsymbol{\beta} + \boldsymbol{\varepsilon}$ mit $\mathbf{X}$ als $n \times (p+1)$ Design-Matrix (inkl. Eins-Spalte für Intercept). \textbf{OLS-Lösung:}
$$\hat{\boldsymbol{\beta}} = (\mathbf{X}^T\mathbf{X})^{-1}\mathbf{X}^T\mathbf{y} \quad \text{(falls } \mathbf{X}^T\mathbf{X} \text{ invertierbar)}$$}

\textbf{Gütemaße:} $R^2$ steigt mit jedem zusätzlichen Prädiktor (auch unwichtigem); \textbf{Adjusted} $R^2_{\text{adj}} = 1 - \frac{(1-R^2)(n-1)}{n-p-1}$ bestraft Modellkomplexität; RSE $= \sqrt{\frac{\text{RSS}}{n-p-1}}$ schätzt $\sigma$.

\textbf{Inferenz:} $t$-Test für einzelne Koeffizienten ($H_0: \beta_j = 0$); $F$-Test für Gesamtmodell.

\subsubsection*{Vergleich: SLR / MLR / Multivariat}

\begin{center}
\small
\begin{tabular}{|l|c|c|}
\hline
\textbf{Modell} & \textbf{Prädiktoren} & \textbf{Response} \\
\hline
Einfach (SLR) & 1 & 1 \\
Multipel (MLR) & $p > 1$ & 1 \\
Multivariat & $p \geq 1$ & $k > 1$ (Koeffizientenmatrix $\mathbf{B}$) \\
\hline
\end{tabular}
\end{center}

\subsubsection*{Logistische Regression}

\definition{\textbf{Zweck:} Regression für binäre Response $y \in \{0,1\}$. Lineare Regression ungeeignet (Vorhersagen außerhalb $[0,1]$). Logistische Regression modelliert $p_i = P(y_i=1 \mid \mathbf{x}_i)$ via Sigmoid:
$$p_i = \frac{1}{1 + e^{-(\beta_0 + \boldsymbol{\beta}^T\mathbf{x}_i)}}$$
\textbf{Log-Odds (Logit):} $\log\!\left(\frac{p_i}{1-p_i}\right) = \beta_0 + \beta_1 x_{i1} + \cdots + \beta_p x_{ip}$ -- linear in den Prädiktoren.}

\textbf{Koeffizienteninterpretation:} Anstieg von $x_j$ um 1 Einheit multipliziert die Odds mit $e^{\beta_j}$.

\textbf{Schätzung:} Maximum Likelihood (nicht OLS), iterativ (z.\,B. Newton-Raphson).

\textbf{Gütemaße:} Deviance; McFadden-$R^2$ (Pseudo-$R^2$); Confusion Matrix (Accuracy, Sensitivity, Specificity); ROC-Kurve und AUC.


% ============================================================

\subsection{Polynomregression, Splines und additive Modelle}
% ============================================================

\definition{\textbf{Polynomregression:} Erweiterung der linearen Regression auf polynomiale Funktionen:
$$y_i = \beta_0 + \beta_1 x_i + \beta_2 x_i^2 + \cdots + \beta_d x_i^d + \varepsilon_i$$
Durch Substitution $x_{ij} = x_i^j$ ergibt sich ein Standard-MLR-Modell, geschätzt via OLS. \textbf{Problem:} Höhere Grade $d$ führen zu Overfitting (hohe Varianz, schlechte Generalisierung). Wahl von $d$ via Cross-Validation, AIC oder BIC. \textbf{Nachteil:} Globale Struktur -- lokal starke Abweichungen erzwingen hohe Polynomgrade überall.}

\definition{\textbf{Splines (Cubic Spline):} Stückweise polynomiale Funktionen. Der $x$-Bereich wird an $K$ \textbf{Knots} in Intervalle aufgeteilt; in jedem Intervall wird ein kubisches Polynom gefittet, mit Stetigkeits- und Glattheits\-bedingungen (Funktion, erste und zweite Ableitung stetig) an den Knots. \textbf{Natural Cubic Spline:} Außerhalb der äußeren Knots ist die Funktion linear -- stabilisiert Extrapolation. Als Linearkombination von Basis-Funktionen darstellbar:
$$f(x) = \sum_{j=1}^M \beta_j B_j(x)$$
Das Modell ist linear in den Koeffizienten und wird mit OLS geschätzt. Anzahl Freiheitsgrade $\approx K + 3$.}

\definition{\textbf{Smoothing Splines:} Regularisierter Ansatz -- minimiere:
$$\sum_{i=1}^n (y_i - f(x_i))^2 + \lambda \int [f''(x)]^2\,dx$$
Erster Term: Datenanpassung; zweiter Term: Strafterm für Rauheit. $\lambda = 0$: perfekte Interpolation; $\lambda \to \infty$: lineare Regressionsgerade. Optimales $\lambda$ via Cross-Validation (GCV oder LOOCV).}

\begin{center}
\small
\begin{tabular}{|l|l|l|}
\hline
\textbf{Aspekt} & \textbf{Polynomregression} & \textbf{Splines} \\
\hline
Flexibilität & Global (abhängig von $d$) & Lokal (variabel in Intervallen) \\
\hline
Overfitting-Risiko & Höher bei großem $d$ & Besser kontrollierbar \\
\hline
Glattheit & Automatisch Grad $d-1$ & Kontrollierbar (typisch $C^2$) \\
\hline
Extrapolation & Polynomial explodiert & Oft linear außerhalb \\
\hline
\end{tabular}
\end{center}

\definition{\textbf{Generalized Additive Models (GAM):} Flexibles Regressionsframework mit additiver Struktur:
$$y_i = \beta_0 + f_1(x_{i1}) + f_2(x_{i2}) + \cdots + f_p(x_{ip}) + \varepsilon_i$$
Jede $f_j$ ist eine glatte Funktion (z.\,B. Smoothing Spline). \textbf{Schätzung via Backfitting:} Iterativ -- für jede Variable $j$ berechne Residuen abzüglich aller anderen $f_k$, glätte diese gegen $x_j$; bis Konvergenz. Ergebnis ist wieder ein regularisiertes lineares Problem. \textbf{Variablen-Auswahl:} Via Partial Dependence Plots, ANOVA-Tests oder Regularisierung. Kompromiss zwischen Interpretierbarkeit (Partial Plots nötig) und Flexibilität.}

\textbf{Praktische Wahl:} Polynomregression bei vermuteter quadratischer/kubischer Beziehung; Splines für glatte nichtlineare Beziehungen mit lokal variierender Komplexität; GAM bei vielen Prädiktoren mit komplexen Abhängigkeiten.

% ============================================================
\newpage

\subsection{PCA, PCR und MDS}
% ============================================================

\definition{\textbf{Principal Component Analysis (PCA):} Unsupervisierte Dimensionsreduktion. Gegeben Datenmatrix $\mathbf{X}$ ($n \times p$); gesucht neue orthogonale Koordinaten (\emph{Principal Components}), in denen maximale Varianz in den ersten Achsen konzentriert ist. PCA-Zerlegung:
$$\mathbf{X} = \mathbf{T}\mathbf{P}^T + \mathbf{E}$$
wobei $\mathbf{T}$ ($n \times a$) die \emph{Score-Matrix} (Koordinaten der Beobachtungen) und $\mathbf{P}$ ($p \times a$) die \emph{Loading-Matrix} (Eigenvektoren der Kovarianzmatrix) ist. Scores sind unkorreliert: $\mathbf{T}^T\mathbf{T} = \text{diag}(\lambda_1,\ldots,\lambda_a)$.}

\textbf{Algorithmus:} (1) Standardisierung von $\mathbf{X}$ (Mittelwert 0, SD 1); (2) Kovarianzmatrix $\mathbf{C} = \frac{1}{n-1}\mathbf{X}^T\mathbf{X}$; (3) Eigenwertzerlegung: $\lambda_1 \geq \lambda_2 \geq \cdots \geq \lambda_p \geq 0$; (4) Scores: $\mathbf{T} = \mathbf{X}\mathbf{P}$.

\textbf{Varianzzerlegung:} $\lambda_j$ = Varianz der $j$-ten PC. Kumulativer Varianzanteil: $\frac{\sum_{j=1}^a \lambda_j}{\sum_{j=1}^p \lambda_j}$. \emph{Scree-Plot}: Eigenwerte gegen PC-Nummer, Wahl von $a$ am „Ellbogen".

\textbf{Visualisierung:} \emph{Score-Plot} (Beobachtungen in PC1/PC2, zeigt Cluster/Ausreißer); \emph{Loading-Plot} (Beitrag der Originalvariablen zu jeder PC); \emph{Biplot} (kombiniert beides).

\definition{\textbf{Principal Component Regression (PCR):} Löst Multikollinearität und $p > n$ in MLR. Statt Originalprädiktoren $\mathbf{X}$ werden die ersten $a$ PCA-Scores als Prädiktoren genutzt:
$$\mathbf{y} = \mathbf{T}\boldsymbol{\vartheta} + \boldsymbol{\varepsilon}, \qquad \hat{\boldsymbol{\vartheta}} = (\mathbf{T}^T\mathbf{T})^{-1}\mathbf{T}^T\mathbf{y}$$
Scores sind orthogonal $\Rightarrow$ OLS stabil und wohldefiniert. Rücktransformation: $\hat{\boldsymbol{\beta}}_{\text{PCR}} = \mathbf{P}\hat{\boldsymbol{\vartheta}}$. Wahl von $a$ via Cross-Validation.}

\definition{\textbf{Multidimensional Scaling (MDS):} Dimensionsreduktion aus einer \textbf{Distanzmatrix} $\mathbf{D}$ ($n \times n$, kein Rohdatensatz nötig). Ziel: Finde niedrigdimensionale Koordinaten $\mathbf{Z}$ ($n \times a$), sodass euklidische Abstände in $\mathbf{Z}$ den ursprünglichen Abständen $d_{ij}$ möglichst entsprechen.

\textbf{Varianten:} \emph{Metric MDS} reproduziert absolute Distanzen (Stressminimierung: $\sqrt{\frac{\sum(d_{ij}-\hat{d}_{ij})^2}{\sum d_{ij}^2}}$); \emph{Non-Metric MDS} erhält nur die Ordnung der Distanzen (robuster gegen Fehler). \textbf{Klassisches MDS:} Double Centering von $\mathbf{D}$ $\to$ Gram-Matrix $\to$ Eigenwertzerlegung $\to$ Koordinaten.}

\begin{center}
\small
\begin{tabular}{|l|l|l|}
\hline
\textbf{Aspekt} & \textbf{PCA} & \textbf{MDS} \\
\hline
Input & Datenmatrix $\mathbf{X}$ & Distanzmatrix $\mathbf{D}$ \\
\hline
Ziel & Maximiere Varianz & Reproduziere Distanzen \\
\hline
Methode & Eigenwertzerlegung Kovarianz & Eigenwertzerlegung Gram-Matrix \\
\hline
Interpretation & Varianz-Richtungen & Objekt-Ähnlichkeiten \\
\hline
\end{tabular}
\end{center}

\textbf{Anwendungen:} PCA: Bioinformatik, Bildkompression, Feature Engineering. PCR: Spektroskopie, Genexpression ($p \gg n$). MDS: Psychometrie, genetische Distanzen, chemische Ähnlichkeitsmaße.

% ============================================================
\newpage

\subsection{Lasso und Ridge-Regression}
% ============================================================

\definition{\textbf{Regularisierung:} Techniken zur Begrenzung der Koeffizientengröße, um Overfitting bei Multikollinearität oder $p \gg n$ zu vermeiden. OLS wird instabil, wenn Eigenwerte von $\mathbf{X}^T\mathbf{X}$ nahe Null sind.}

\definition{\textbf{Ridge Regression (L2-Penalty):} Minimiere:
$$L_{\text{ridge}} = \|\mathbf{y} - \mathbf{X}\boldsymbol{\beta}\|_2^2 + \lambda\|\boldsymbol{\beta}\|_2^2 = \text{RSS} + \lambda\sum_j \beta_j^2$$
Geschlossene Lösung: $\hat{\boldsymbol{\beta}}_{\text{ridge}} = (\mathbf{X}^T\mathbf{X} + \lambda\mathbf{I})^{-1}\mathbf{X}^T\mathbf{y}$. Für orthonormale Prädiktoren: $\hat{\boldsymbol{\beta}}_{\text{ridge}} = \frac{1}{1+\lambda}\hat{\boldsymbol{\beta}}_{\text{OLS}}$ -- alle Koeffizienten werden gleichmäßig geschrumpft, aber \emph{nie exakt zu Null}.}

\definition{\textbf{Lasso Regression (L1-Penalty):} Minimiere:
$$L_{\text{lasso}} = \|\mathbf{y} - \mathbf{X}\boldsymbol{\beta}\|_2^2 + \lambda\|\boldsymbol{\beta}\|_1 = \text{RSS} + \lambda\sum_j |\beta_j|$$
Keine geschlossene Lösung; numerisch via Coordinate Descent. \textbf{Kritischer Unterschied:} L1-Penalty setzt Koeffizienten \emph{exakt zu Null} $\to$ automatische Variablen-Auswahl. \textbf{Geometrische Intuition:} Lasso-Constraint ist ein Diamant (spitze Ecken auf den Achsen); Ridge ist ein Kreis. Die Ecken des Diamanten machen Nulllösungen wahrscheinlich.}

\begin{center}
\small
\begin{tabular}{|l|l|l|}
\hline
\textbf{Aspekt} & \textbf{Ridge (L2)} & \textbf{Lasso (L1)} \\
\hline
Penalty & $\lambda\sum\beta_j^2$ & $\lambda\sum|\beta_j|$ \\
\hline
Lösung & Geschlossen & Numerisch (Coord. Descent) \\
\hline
Koeff. $= 0$? & Selten & Oft (Variablen-Auswahl) \\
\hline
Multikollinearität & Gut behandelt & Wählt eine Variable aus \\
\hline
Ideal für & Alle Vars. relevant & Sparse, $p \gg n$ \\
\hline
\end{tabular}
\end{center}

\definition{\textbf{Elastic Net:} Kombiniert beide Penalties mit Mischungsparameter $\alpha \in [0,1]$:
$$L_{\text{EN}} = \text{RSS} + \lambda\bigl[\alpha\|\boldsymbol{\beta}\|_1 + (1-\alpha)\|\boldsymbol{\beta}\|_2^2\bigr]$$
$\alpha=0$: Ridge; $\alpha=1$: Lasso; dazwischen: Kompromiss aus Variablen-Auswahl und Stabili\-tät bei korrelierten Prädiktoren.}

\textbf{Bias-Varianz Tradeoff:} Kleines $\lambda$: niedriger Bias, hohe Varianz (nahe OLS); großes $\lambda$: hoher Bias, niedrige Varianz. Optimales $\lambda$ (und $\alpha$ bei Elastic Net) immer via \textbf{Cross-Validation}.

\textbf{Solution Path:} Visualization, wie Koeffizienten sich bei variierendem $\lambda$ verändern. Ridge: kontinuierliche Schrumpfung; Lasso: sprunghafte Nullsetzung einzelner Variablen.

% ============================================================
\newpage

\subsection{Entscheidungsbäume und Random Forests}
% ============================================================

\definition{\textbf{Decision Tree:} Supervisierter Algorithmus für Klassifikation und Regression. Grundstruktur: \emph{Knoten} (Bedingung $x_j \leq c$), \emph{Äste} (ja/nein), \emph{Blätter} (Klassen- oder Vorhersagewert). Neue Beobachtungen traversieren den Baum von der Wurzel bis zu einem Blatt.}

\definition{\textbf{Baumaufbau -- Recursive Partitioning (Top-Down Greedy):}
\begin{enumerate}[leftmargin=1.5em]
    \item Alle Trainingsdaten an der Wurzel.
    \item Finde Variable $x_j$ und Schwellwert $c$, der die reinsten Kindknoten liefert.
    \item Teile in $x_j \leq c$ (links) und $x_j > c$ (rechts); rekursiv wiederholen.
    \item Stopp bei: alle Punkte einer Klasse, minimaler Knotengröße, maximaler Tiefe oder keiner weiteren Verbesserung.
\end{enumerate}}

\definition{\textbf{Unreinheitsmaße:}

\emph{Klassifikation:} \textbf{Gini-Index} $\text{Gini} = 1 - \sum_k p_k^2$ (0 = rein); \textbf{Entropie} $\text{Entropy} = -\sum_k p_k\log_2 p_k$; \textbf{Information Gain} $= \text{Entropy(Parent)} - \sum_i \frac{n_i}{n}\text{Entropy(Child}_i)$.

\emph{Regression:} Minimierung der kombinierten RSS $= \sum_i (y_i - \hat{y})^2$ in beiden Kindknoten.}

\definition{\textbf{Overfitting und Pruning:} Tiefe Bäume fitten Trainingsdaten perfekt, generalisieren aber schlecht. Gegenmaßnahmen: \emph{Early Stopping} (Tiefenlimit, Mindestgröße pro Knoten); \emph{Post-Pruning} (vollen Baum trainieren, dann von unten Knoten entfernen wenn Validierungsperformance steigt); \emph{Cost Complexity Pruning} (minimiert Trainingsfehler + Komplexitätsstrafe).}

\begin{center}
\small
\begin{tabular}{|l|l|}
\hline
\textbf{Vorteile} & \textbf{Nachteile} \\
\hline
Interpretierbar und visualisierbar & Instabil (kleine Datenänderungen) \\
Keine Skalierung nötig & Overfitting ohne Regularisierung \\
Kategoriale + numerische Prädiktoren & Greedy, nicht global optimal \\
Erfasst Nichtlinearität \& Interaktionen & Bias bei unbalancierten Klassen \\
\hline
\end{tabular}
\end{center}

\definition{\textbf{Random Forest:} Ensemble aus $B$ Entscheidungsbäumen mit zwei Schlüsselmechanismen: \textbf{Bootstrap-Sampling} (jeder Baum auf Zufallsstichprobe mit Zurücklegen) und \textbf{Random Feature Selection} (bei jedem Split nur $m \approx \sqrt{p}$ zufällig gewählte Prädiktoren). Beides reduziert Varianz und Korrelation zwischen Bäumen. Vorhersage via Mehrheitsvote (Klassifikation) oder Durchschnitt (Regression).}

\definition{\textbf{Out-of-Bag (OOB) Error:} Ca.\ 37\% der Trainingsdaten sind in jeder Bootstrap-Stichprobe nicht enthalten. Diese können zur Validierung genutzt werden -- ohne separaten Validierungssatz nötig.}

\definition{\textbf{Feature Importance:} \emph{Mean Decrease in Impurity (MDI)}: Summiere Impurity-Reduktion bei Splits auf $X_j$ über alle Bäume. \emph{Permutation Importance (MDA)}: Permutiere $X_j$ zufällig, messe OOB-Fehleranstieg; stärkere Verschlechterung $=$ wichtigere Variable.}

\begin{center}
\small
\begin{tabular}{|l|l|l|}
\hline
\textbf{Aspekt} & \textbf{Decision Tree} & \textbf{Random Forest} \\
\hline
Varianz & Hoch & Niedrig (Ensemble) \\
Interpretierbarkeit & Hoch & Niedrig (Black Box) \\
Performance & Moderat & Höher \\
Feature Importance & Einfach & MDI oder Permutation \\
\hline
\end{tabular}
\end{center}

\definition{\textbf{Weitere Ensemble-Methoden:} \emph{Bagging}: Bootstrap-Aggregation auf beliebige Lerner. \emph{Boosting} (AdaBoost, Gradient Boosting): Bäume sequenziell trainieren, jeder korrigiert Fehler des vorigen -- höhere Accuracy, aber overfitting-anfälliger. \emph{XGBoost}: Regularisierte, hochoptimierte GBM-Implementierung, weit verbreitet in der Praxis.}

% ============================================================
\newpage

\subsection{Linear Discriminant Analysis (LDA) und Quadratic Discriminant Analysis (QDA)}
% ============================================================

\definition{\textbf{Klassifikations-Setting:} Datenmatrix $\mathbf{X}$ ($n \times p$), kategoriale Response $y \in \{1,\ldots,K\}$. Beobachtungen in Klasse $k$ entstammen Dichte $f_k(\mathbf{x})$. \textbf{Bayes-Klassifikationsregel} mit Prior-Wahrscheinlichkeiten $p_k$:
$$\hat{k} = \arg\max_{k=1,\ldots,K}\; f_k(\mathbf{x}) \cdot p_k$$
Intuition: Selbst wenn $f_2(\mathbf{x}) > f_1(\mathbf{x})$, kann Klasse 1 bevorzugt werden, wenn $p_1 \gg p_2$.}

\definition{\textbf{LDA-Annahmen:} Multivariat normalverteilt innerhalb jeder Klasse, $\mathbf{x}|C{=}k \sim N_p(\boldsymbol{\mu}_k, \boldsymbol{\Sigma})$, mit \emph{gemeinsamer} Kovarianzmatrix $\boldsymbol{\Sigma}$ für alle Klassen $\Rightarrow$ \textbf{lineare} Entscheidungsgrenzen. \textbf{Diskriminanzfunktion:}
$$\delta_k(\mathbf{x}) = \mathbf{x}^T\boldsymbol{\Sigma}^{-1}\boldsymbol{\mu}_k - \tfrac{1}{2}\boldsymbol{\mu}_k^T\boldsymbol{\Sigma}^{-1}\boldsymbol{\mu}_k + \log p_k$$
Klassifikation: $\hat{k} = \arg\max_k \delta_k(\mathbf{x})$. Entscheidungsgrenzen zwischen zwei Klassen ($\delta_k = \delta_{k'}$) sind lineare Hyperebenen.}

\textbf{Parameterschätzung:} Klassenmittelwerte $\hat{\boldsymbol{\mu}}_k = \frac{1}{n_k}\sum_{i:y_i=k}\mathbf{x}_i$; Priors $\hat{p}_k = n_k/n$; gepoolte Kovarianzmatrix $\hat{\boldsymbol{\Sigma}}_{\text{pooled}} = \frac{1}{n-K}\sum_k (n_k-1)\hat{\boldsymbol{\Sigma}}_k$.

\definition{\textbf{QDA-Annahmen:} Wie LDA, aber jede Klasse hat eine \emph{eigene} Kovarianzmatrix $\boldsymbol{\Sigma}_k$ $\Rightarrow$ \textbf{quadratische} Entscheidungsgrenzen (Quadriken). \textbf{Diskriminanzfunktion:}
$$\delta_k(\mathbf{x}) = -\tfrac{1}{2}\log|\boldsymbol{\Sigma}_k| - \tfrac{1}{2}(\mathbf{x}-\boldsymbol{\mu}_k)^T\boldsymbol{\Sigma}_k^{-1}(\mathbf{x}-\boldsymbol{\mu}_k) + \log p_k$$
Der quadratische Term $\mathbf{x}^T\boldsymbol{\Sigma}_k^{-1}\mathbf{x}$ unterscheidet QDA von LDA. Pro Klasse werden $p(p+1)/2$ Kovarianzparameter geschätzt -- erhöhtes Overfitting-Risiko bei kleinem $n$.}

\begin{center}
\small
\begin{tabular}{|l|l|l|}
\hline
\textbf{Aspekt} & \textbf{LDA} & \textbf{QDA} \\
\hline
Kovarianzannahme & Gemeinsam $\boldsymbol{\Sigma}$ & Klassenspezifisch $\boldsymbol{\Sigma}_k$ \\
\hline
Entscheidungsgrenze & Linear & Quadratisch \\
\hline
Parameteranzahl & Wenige & Viele (pro Klasse) \\
\hline
Bias/Varianz & Höherer Bias & Höhere Varianz \\
\hline
Bevorzugt wenn & Ähnliche Spreads, kleines $n$ & Unterschiedliche Spreads, großes $n$ \\
\hline
\end{tabular}
\end{center}

\definition{\textbf{Varianten:}
\begin{itemize}[leftmargin=1.5em]
    \item \textbf{Fisher's Linear Discriminant} (binär): Findet Projektionsrichtung $\mathbf{w} \propto \boldsymbol{\Sigma}_{\text{pooled}}^{-1}(\bar{\mathbf{x}}_1 - \bar{\mathbf{x}}_2)$, die Between-Class- zu Within-Class-Variabilität maximiert. Konzeptuell auf Dimensionsreduktion ausgerichtet, ohne Normalverteilungsannahme.
    \item \textbf{Regularized DA (RDA):} Interpoliert via $\hat{\boldsymbol{\Sigma}}_k(\lambda) = \lambda\hat{\boldsymbol{\Sigma}}_{\text{pooled}} + (1-\lambda)\hat{\boldsymbol{\Sigma}}_k$, $\lambda \in [0,1]$ per Cross-Validation. Kompromiss zwischen LDA ($\lambda=1$) und QDA ($\lambda=0$).
    \item \textbf{Mixture DA (MDA):} Modelliert jede Klasse als Mischung mehrerer Normalverteilungen -- für multimodale Klassenverteilungen.
\end{itemize}}

\textbf{Methodenwahl:} LDA bei ähnlichen Kovarianzstrukturen, kleinem $n$; QDA bei deutlich unterschiedlichen Spreads, großem $n$ ($n \gg p \cdot K$); RDA bei Unsicherheit; andere Klassifikatoren (SVM, kNN, Trees) bei stark verletzten Normalverteilungsannahmen.

% ============================================================
\newpage

\subsection{k-Nearest Neighbors (kNN)}
% ============================================================

\definition{\textbf{kNN-Grundkonzept:} Nicht-parametrischer \emph{Lazy Learner} -- keine Trainingsphase, Trainingsdaten werden direkt gespeichert. Für neue Beobachtung $\mathbf{x}_0$: (1) Berechne Distanzen zu allen $n$ Trainingspunkten; (2) identifiziere die $k$ nächsten Nachbarn; (3) \textbf{Klassifikation}: Mehrheitsvote; \textbf{Regression}: Durchschnitt der $k$ Zielwerte.}

\definition{\textbf{Distanzmaße:} Eine gültige Distanzfunktion muss vier Metrik-Axiome erfüllen: Nicht-Negativität, Identität, Symmetrie, Dreiecksungleichung. Gebräuchliche Metriken:
\begin{itemize}[leftmargin=1.5em]
    \item \textbf{Minkowski-Distanz (Verallgemeinerung):} $d_r = \bigl(\sum_j |x_{A,j}-x_{B,j}|^r\bigr)^{1/r}$; mit $r=1$: Manhattan, $r=2$: Euklidisch, $r\to\infty$: Chebyshev ($\max_j|x_{A,j}-x_{B,j}|$).
    \item \textbf{Mahalanobis-Distanz:} $\sqrt{(\mathbf{x}_A-\mathbf{x}_B)^T\boldsymbol{\Sigma}^{-1}(\mathbf{x}_A-\mathbf{x}_B)}$ -- berücksichtigt Korrelation und Skalierung.
    \item \textbf{Tanimoto-Distanz (binäre Vektoren):} $d = 1 - \frac{\mathbf{x}_A^T\mathbf{x}_B}{\|\mathbf{x}_A\|_1 + \|\mathbf{x}_B\|_1 - \mathbf{x}_A^T\mathbf{x}_B}$
\end{itemize}
\textbf{Wichtig:} Variablen vor kNN immer normalisieren/standardisieren -- unterschiedliche Skalen dominieren sonst die Distanzberechnung.}

\definition{\textbf{Einfluss von $k$ -- Bias-Varianz Tradeoff:}
\begin{itemize}[leftmargin=1.5em]
    \item \textbf{Kleines $k$ (z.\,B. $k=1$):} Flexible, zackige Entscheidungsgrenzen; hohes Overfitting-Risiko (hohe Varianz).
    \item \textbf{Großes $k$ (z.\,B. $k=n$):} Glatte, einfache Grenzen; hohes Underfitting-Risiko (hoher Bias).
\end{itemize}
Optimales $k$ via Cross-Validation; bei binärer Klassifikation ungerades $k$ empfohlen (kein Unentschieden).}

\begin{center}
\small
\begin{tabular}{|l|l|}
\hline
\textbf{Vorteile} & \textbf{Nachteile} \\
\hline
Einfach, keine Trainingsphase & $O(n \cdot p)$ Aufwand pro Vorhersage \\
Nicht-parametrisch, flexibel & Sensibel gegenüber Skalierung \\
Klassifikation und Regression & Curse of Dimensionality bei großem $p$ \\
Keine Verteilungsannahmen & Ausreißer nah am Testpunkt schädlich \\
\hline
\end{tabular}
\end{center}

\definition{\textbf{Varianten:} \emph{Gewichtetes kNN}: Gewichte Nachbarn nach Distanz ($w_i = 1/d_i$ oder Gaußkern) -- nähere Nachbarn haben mehr Einfluss. \emph{Approximate kNN}: KD-Trees oder Ball-Trees zur Beschleunigung bei großen $n$. \emph{Feature Selection}: PCA oder Relevanzgewichtung verbessert kNN bei hochdimensionalen Daten.}

\textbf{Anwendungen:} Klassifikation bei kleinem $p$ (Iris, Handschriften); Anomalieerkennung (Punkte weit von allen anderen); Recommender Systems (ähnliche Kunden/Produkte).











% ============================================================

% ============================================================

\subsection{Ähnlichkeitsmaße und Distanzmaße}
% ============================================================

\definition{\textbf{Grundkonzept:} Distanz $d$ misst Unähnlichkeit; Ähnlichkeit $s$ ist deren Inverses. Bei normierten Maßen gilt $s = 1 - d$. Objekte werden als $p$-dimensionale Vektoren $\mathbf{x}_A, \mathbf{x}_B \in \mathbb{R}^p$ dargestellt. Die \textbf{Distanzmatrix} $\mathbf{D} = (d_{ij})$ ist eine $n\times n$-Matrix mit $d_{ii}=0$, $d_{ij} \geq 0$ und $d_{ij} = d_{ji}$ (symmetrisch); sie enthält $\frac{n(n-1)}{2}$ nicht-redundante Einträge.}

\noindent\textbf{Distanzmaße für numerische (metrische) Daten:}

\begin{itemize}[leftmargin=1.5em]
    \item \textbf{Euklidische Distanz:}
    $d_{\text{Euk}} = \sqrt{\sum_{j=1}^{p}(x_{B,j}-x_{A,j})^2} = \|\mathbf{x}_B - \mathbf{x}_A\|_2$
    -- geometrischer Luftlinienabstand; empfindlich gegenüber Ausreißern und unterschiedlichen Skalierungen (Standardisierung empfohlen).

    \item \textbf{Manhattan-Distanz (City Block):}
    $d_{\text{Man}} = \sum_{j=1}^{p}|x_{B,j}-x_{A,j}|$
    -- Summe absoluter Koordinatendifferenzen; robuster gegenüber Ausreißern.

    \item \textbf{Minkowski-Distanz:}
    $d_{p'} = \bigl(\sum_{j=1}^{p}|x_{B,j}-x_{A,j}|^{p'}\bigr)^{1/p'}$
    -- Verallgemeinerung; $p'=1$: Manhattan; $p'=2$: Euklidisch; $p'\to\infty$: Chebyshev-Distanz $\max_j|x_{A,j}-x_{B,j}|$ (größte Koordinatendifferenz).

    \item \textbf{Mahalanobis-Distanz:}
    $d_{\text{Mah}} = \sqrt{(\mathbf{x}_B-\mathbf{x}_A)^T\hat{\boldsymbol{\Sigma}}^{-1}(\mathbf{x}_B-\mathbf{x}_A)}$
    -- berücksichtigt Kovarianzstruktur und Skalierungen; keine separate Standardisierung nötig. Punkte gleicher Mahalanobis-Distanz bilden Ellipsen (vs. Kreise bei Euklidisch). Nachteil: erfordert invertierbare $\hat{\boldsymbol{\Sigma}}$, instabil bei $p > n$.
\end{itemize}

\noindent\textbf{Ähnlichkeitsmaße:}

\begin{itemize}[leftmargin=1.5em]
    \item \textbf{Kosinus-Ähnlichkeit:}
    $\cos(\omega) = \frac{\mathbf{x}_A^T\mathbf{x}_B}{\|\mathbf{x}_A\|_2\|\mathbf{x}_B\|_2} \in [-1,1]$
    -- misst den Winkel zwischen Vektoren, \emph{unabhängig von ihrer Länge}; Standard für Text- und Spektrendaten (TF-IDF, NIR).

    \item \textbf{Tanimoto-Koeffizient (Jaccard-Index)} -- für \textbf{Binärvektoren} $\mathbf{x}_A, \mathbf{x}_B \in \{0,1\}^p$:
    $$t_{AB} = \frac{\sum_j \text{AND}(x_{A,j}, x_{B,j})}{\sum_j \text{OR}(x_{A,j}, x_{B,j})} = \frac{\mathbf{x}_A^T\mathbf{x}_B}{\mathbf{x}_A^T\mathbf{1} + \mathbf{x}_B^T\mathbf{1} - \mathbf{x}_A^T\mathbf{x}_B} \in [0,1]$$
    Zähler: Anzahl gemeinsamer Einsen (AND); Nenner: Anzahl Positionen mit mindestens einer Eins (OR). Abgeleitete Distanz: $d_{\text{Tan}} = 1 - t_{AB} = \frac{\sum_j \text{XOR}(x_{A,j},x_{B,j})}{\sum_j \text{OR}(x_{A,j},x_{B,j})}$. Anwendung: Chemoinformatik (Molekülstrukturen), Genetik.

    \item \textbf{Dice-Koeffizient:}
    $S_{\text{Dice}} = \frac{2|A \cap B|}{|A| + |B|} = \frac{2\mathbf{x}_A^T\mathbf{x}_B}{\mathbf{x}_A^T\mathbf{1} + \mathbf{x}_B^T\mathbf{1}}$
    -- gewichtet gemeinsame Einsen stärker als Tanimoto. Zusammenhang: $S_{\text{Dice}} = \frac{2t_{AB}}{1+t_{AB}}$. Einsatz in NLP und Biologie.
\end{itemize}

\noindent\textbf{Korrelationsbasierte Distanzmaße:} Für $p$-dimensionale Vektoren $\mathbf{x}_A$, $\mathbf{x}_B$ (Beobachtungen als Vektoren über Variablen):
\begin{itemize}[leftmargin=1.5em]
    \item \textbf{Pearson-Distanz:} $d = 1 - r_{\text{Pearson}}$ -- misst lineare Abhängigkeit; invariant gegenüber linearen Transformationen.
    \item \textbf{Spearman-Distanz:} $d = 1 - r_{\text{Spearman}}$ -- rangbasiert, robust gegenüber Ausreißern und monoton nichtlinearen Zusammenhängen.
    \item \textbf{Kendall-Distanz:} $d = 1 - r_{\text{Kendall}}$ -- basiert auf konkordanten/diskordanten Paaren.
\end{itemize}

\noindent\textbf{Distanzmaße für ordinale und nominale Variablen:}

Für \textbf{ordinale} Variablen mit $m$ geordneten Stufen werden die Ausprägungen durch $\frac{i-1/2}{m}$ ($i=1,\ldots,m$) in numerische Werte transformiert, danach können metrische Distanzmaße angewendet werden. Für \textbf{nominale} Variablen wird eine Distanzmatrix $\mathbf{M}$ mit $m_{ii}=0$ und $m_{ii'}=1$ für $i\neq i'$ (einfachster Fall: 0-1-Kodierung) explizit definiert.

Bei \textbf{gemischten Datentypen} (quantitativ und qualitativ) wird die Gesamtdistanz als gewichtetes Mittel variablenweiser Distanzen berechnet:
$$D(\mathbf{x}_A, \mathbf{x}_B) = \sum_{j=1}^{p} w_j \cdot d_j(x_{A,j}, x_{B,j}), \quad \sum_j w_j = 1$$
\textbf{Wichtig:} Unterschiedliche Skalierungen der Variablen führen dazu, dass Variablen mit großem Wertebereich die Gesamtdistanz dominieren -- \textbf{Standardisierung} ist daher wesentlich, um jedem Merkmal gleichen Einfluss zu geben.

\begin{center}
\small
\begin{tabular}{|l|l|l|}
\hline
\textbf{Situation} & \textbf{Empfohlenes Maß} & \textbf{Skalenniveau} \\
\hline
Normalisierte numerische Daten & Euklidisch & Metrisch \\
\hline
Korrelierte Variablen & Mahalanobis & Metrisch \\
\hline
Längenunabhängige Vektoren & Kosinus & Metrisch \\
\hline
Binäre Merkmalsvektoren & Tanimoto / Dice & Nominal (binär) \\
\hline
Token-Überlappung (NLP) & Dice & Nominal (binär) \\
\hline
Rangbasierte / robuste Analyse & Spearman-Distanz & Ordinal / metrisch \\
\hline
Ausreißer vorhanden & Manhattan oder Mahalanobis & Metrisch \\
\hline
Ordinale Kategorien & Transformation + Euklidisch & Ordinal \\
\hline
Nominale Kategorien & 0-1-Distanzmatrix & Nominal \\
\hline
\end{tabular}
\end{center}

\subsection{K-means, K-medoids und Fuzzy-Clustering}
% ============================================================

\definition{\textbf{K-means:} Partitioniert $n$ Beobachtungen in $K$ disjunkte Cluster durch Minimierung der Within-Cluster-Sum-of-Squares (WCSS):
$$W = \sum_{k=1}^{K}\sum_{i \in C_k} \|\mathbf{x}_i - \mathbf{c}_k\|_2^2$$
wobei $\mathbf{c}_k = \frac{1}{n_k}\sum_{i \in C_k}\mathbf{x}_i$ das Zentroid von Cluster $k$ ist. $K$ muss vorab festgelegt werden.}

\textbf{Algorithmus:}
\begin{enumerate}[leftmargin=1.5em]
    \item Wähle $K$ zufällige Startpunkte als Zentroide.
    \item \emph{Zuordnung:} Weise jede Beobachtung dem nächsten Zentroid zu.
    \item \emph{Update:} Berechne neue Zentroide als Clustermittelwert.
    \item Wiederhole Schritte 2--3 bis keine Änderungen mehr auftreten.
\end{enumerate}

\definition{\textbf{Wahl von $K$:} Die \textbf{Elbow-Methode} plottet WCSS gegen $K$; der Knick der monoton fallenden Kurve zeigt das sinnvolle $K$. Die \textbf{Hopkins-Statistik}
$$H = \frac{\sum d_U(i)}{\sum d_U(i) + \sum d_W(i)}$$
testet, ob überhaupt eine Clusterstruktur vorliegt: $H \approx 0{,}5$ -- zufällige Verteilung, kein Clustering sinnvoll; $H \to 1$ -- starke Clusterstruktur.}

\definition{\textbf{K-medoids:} Wie K-means, aber Clusterzentren müssen \emph{echte Beobachtungen} (Medoide) sein -- das Medoid minimiert die Summe der Abstände zu allen anderen Clusterpunkten. Vorteile: robuster gegenüber Ausreißern; beliebige Distanzmaße einsetzbar; Zentren sind direkt interpretierbar.}

\definition{\textbf{Fuzzy C-Means (FCM):} Jede Beobachtung $i$ erhält einen Zugehörigkeitsgrad $u_{ij} \in [0,1]$ zu jedem Cluster $j$ (mit $\sum_j u_{ij} = 1$). Minimiert wird:
$$J_m = \sum_{i=1}^{n}\sum_{j=1}^{K} u_{ij}^m \|\mathbf{x}_i - \mathbf{c}_j\|^2$$
Fuzzifikationsparameter $m > 1$ (typisch $m=2$): größeres $m$ $\to$ weichere Zugehörigkeiten. Memberships und Zentroide werden abwechselnd aktualisiert:
$$u_{ij} = \frac{1}{\sum_{l=1}^{K}\!\left(\frac{\|\mathbf{x}_i - \mathbf{c}_j\|}{\|\mathbf{x}_i - \mathbf{c}_l\|}\right)^{\!\frac{2}{m-1}}}, \qquad \mathbf{c}_j = \frac{\sum_i u_{ij}^m \mathbf{x}_i}{\sum_i u_{ij}^m}$$}

\begin{center}
\small
\begin{tabular}{|l|l|l|l|}
\hline
\textbf{Aspekt} & \textbf{K-means} & \textbf{K-medoids} & \textbf{Fuzzy C-Means} \\
\hline
Zuordnung & Hart & Hart & Weich (graduell) \\
\hline
Zentren & Durchschnitt & Echte Beobachtung & Gewichteter Durchschnitt \\
\hline
Ausreißer-Robustheit & Gering & Hoch & Mittel \\
\hline
Distanzmaß & Euklidisch & Beliebig & Euklidisch \\
\hline
Zusatzparameter & -- & -- & Fuzzifikator $m$ \\
\hline
\end{tabular}
\end{center}

% ============================================================
\newpage

\subsection{Model-based Clustering}
% ============================================================

\definition{\textbf{Grundkonzept:} Daten werden als Realisierungen einer Mischung von $K$ Wahrscheinlichkeitsverteilungen modelliert (üblicherweise multivariate Normalverteilungen):
$$f(\mathbf{x}) = \sum_{k=1}^{K} \pi_k\, f_k(\mathbf{x} \mid \boldsymbol{\mu}_k, \boldsymbol{\Sigma}_k)$$
mit Mischungsgewichten $\pi_k > 0$, $\sum_k \pi_k = 1$. Im Gegensatz zu K-means liefert dieser Ansatz ein probabilistisches Framework mit weichen Zuordnungen und statistischer Inferenz.}

Die \textbf{multivariate Normalverteilung} $N_p(\boldsymbol{\mu}, \boldsymbol{\Sigma})$ hat die Dichte:
$$f(\mathbf{x}) = \frac{1}{(2\pi)^{p/2}|\boldsymbol{\Sigma}|^{1/2}} \exp\!\left(-\tfrac{1}{2}(\mathbf{x}-\boldsymbol{\mu})^T\boldsymbol{\Sigma}^{-1}(\mathbf{x}-\boldsymbol{\mu})\right)$$

\definition{\textbf{EM-Algorithmus} -- iterative Maximum-Likelihood-Schätzung:
\begin{enumerate}[leftmargin=1.5em]
    \item \textbf{Initialisierung:} Startparameter (z.\,B. aus K-means).
    \item \textbf{E-Schritt:} Posteriore Zugehörigkeitswahrscheinlichkeit (ähnlich Fuzzy-Memberships):
    $$\tau_{ik} = \frac{\pi_k\, f_k(\mathbf{x}_i \mid \boldsymbol{\theta}_k)}{\sum_{j} \pi_j\, f_j(\mathbf{x}_i \mid \boldsymbol{\theta}_j)}$$
    \item \textbf{M-Schritt:} Aktualisiere $\pi_k$, $\boldsymbol{\mu}_k$ und $\boldsymbol{\Sigma}_k$ als $\tau_{ik}$-gewichtete Mittel.
    \item Wiederhole bis Log-Likelihood konvergiert.
\end{enumerate}}

\textbf{Vorteile:} Probabilistisches Framework; weiche Zuordnungen; Modellwahl via AIC/BIC möglich. \textbf{Nachteile:} Rechenaufwendiger als K-means; Normalverteilungsannahme nicht immer realistisch; Kovarianzmatrizen instabil bei $p \gg n$.

% ============================================================
\newpage

\subsection{Hierarchisches Clustering}
% ============================================================

\definition{\textbf{Grundkonzept:} Erzeugt eine vollständige Hierarchie von Clusterlösungen statt einer einzelnen festen Partition. \textbf{Agglomerative} Verfahren (gebräuchlicher) starten mit $n$ Einzelclustern und fusionieren iterativ die ähnlichsten Paare bis ein Gesamtcluster verbleibt. Divisive Verfahren (umgekehrt) sind rechenintensiver und seltener.}

\definition{\textbf{Linkage-Methoden} definieren den Abstand zwischen zwei Clustern $C_i$ und $C_j$:
\begin{itemize}[leftmargin=1.5em]
    \item \textbf{Complete Linkage:} $\max_{a \in C_i, b \in C_j} d(a,b)$ -- maximaler Abstand; globulare Cluster, kein Chaining.
    \item \textbf{Single Linkage:} $\min_{a \in C_i, b \in C_j} d(a,b)$ -- minimaler Abstand; anfällig für Kettenbildung.
    \item \textbf{Average Linkage (UPGMA):} Mittlerer Abstand aller Punktpaare -- Kompromiss zwischen Complete und Single.
    \item \textbf{Ward's Method:} $\frac{n_i n_j}{n_i+n_j}\|\mathbf{c}_i - \mathbf{c}_j\|^2$ -- minimiert Within-Cluster-Varianz bei Fusion; ähnlich K-means-Kriterium; oft bestes Ergebnis in der Praxis.
\end{itemize}}

\definition{\textbf{Dendrogramm:} Baumdiagramm der Fusionshistorie. Blätter = Einzelbeobachtungen; die Höhe einer Verbindung = Fusionsabstand. Horizontale Schnitte liefern Clusterlösungen für beliebiges $K$.}

\definition{\textbf{Agglomerationskoeffizient (ac):} Misst die Stärke der Clusterstruktur. Für jede Beobachtung $i$: $m(i) = d_{i,\text{first}}/d_{i,\text{last}}$. Mittelwert:
$$ac = \frac{1}{n}\sum_{i=1}^{n}(1 - m(i))$$
Werte nahe 1 deuten auf klare, gut getrennte Cluster hin.}

\textbf{Vorteil:} Kein vorab festgelegtes $K$; Dendrogramm ermöglicht visuelle Exploration; verschiedene Linkage-Methoden erlauben unterschiedliche Clusterformen. \textbf{Nachteil:} $O(n^2)$ Komplexität; Fusionsentscheidungen sind final (keine Neuoptimierung).

% ============================================================
\newpage

\subsection{Bedingte Wahrscheinlichkeit und Bayes-Theorem}
% ============================================================

\definition{\textbf{Bedingte Wahrscheinlichkeit:} Wahrscheinlichkeit von $A$ unter der Bedingung, dass $B$ bereits eingetreten ist:
$$P(A \mid B) = \frac{P(A \cap B)}{P(B)}, \quad P(B) > 0$$
Interpretation: Das Eintreten von $B$ schränkt den Stichprobenraum auf $B$ ein -- $P(A\mid B)$ ist eine Neueinschätzung von $P(A)$ unter dieser Information.}

\textbf{Produktregel (Kettenregel):} $P(A \cap B) = P(A \mid B)\cdot P(B) = P(B \mid A)\cdot P(A)$; für mehrere Ereignisse: $P(A \cap B \cap C) = P(A)\cdot P(B\mid A)\cdot P(C \mid A \cap B)$.

\textbf{Satz der totalen Wahrscheinlichkeit:} Für eine Partition $B_1,\ldots,B_K$ von $\Omega$:
$$P(A) = \sum_{k=1}^{K} P(A \mid B_k)\cdot P(B_k)$$

\definition{\textbf{Satz von Bayes:}
$$P(A \mid B) = \frac{P(B \mid A)\cdot P(A)}{P(B)} = \frac{P(B \mid A)\cdot P(A)}{\sum_k P(B \mid A_k)\cdot P(A_k)}$$
Die vier Terme: \textbf{Posterior} $P(A|B)$ (aktualisiertes Wissen nach Beobachtung), \textbf{Likelihood} $P(B|A)$, \textbf{Prior} $P(A)$ (Vorwissen), \textbf{Evidenz} $P(B)$ (Normalisierungskonstante). Der Satz erlaubt, ``rückwärts'' zu schließen: von bekanntem $P(B|A)$ auf gesuchtes $P(A|B)$.}

\textbf{Beispiel -- Medizinischer Test:} Prävalenz 1\%, Sensitivität 95\%, Spezifität 90\%.
$$P(+) = 0{,}95 \cdot 0{,}01 + 0{,}10 \cdot 0{,}99 = 0{,}1085$$
$$P(\text{krank}\mid +) = \frac{0{,}95 \cdot 0{,}01}{0{,}1085} \approx 8{,}8\%$$
Trotz gutem Test: ein positives Ergebnis bedeutet nur ca.\ 9\% Krankheitswahrscheinlichkeit -- die \textbf{Base-Rate Fallacy} (seltene Erkrankungen erzeugen viele falsch-positive Befunde).

% ============================================================
\newpage

\subsection{Probabilistische Graphische Modelle}
% ============================================================

\definition{\textbf{Graphische Modelle:} Repräsentation von Wahrscheinlichkeitsverteilungen über Graphen, wobei Knoten Zufallsvariablen und Kanten Abhängigkeiten repräsentieren. \textbf{Bayes-Netz:} Ein gerichteter, azyklischer Graph (DAG); eine Kante $A \to B$ drückt aus, dass $A$ direkter Elternknoten von $B$ ist. Die gemeinsame Verteilung faktorisiert über die Graphstruktur:
$$P(X_1, \ldots, X_n) = \prod_{i=1}^{n} P(X_i \mid \text{Eltern}(X_i))$$
Dies reduziert die Parametrisierung erheblich durch bedingte Unabhängigkeiten.}

\definition{\textbf{Bedingte Unabhängigkeit und d-Separation:} Drei kanonische Strukturen bestimmen den Informationsfluss:
\begin{itemize}[leftmargin=1.5em]
    \item \textbf{Seriell} $X \to Y \to Z$: Fluss von $X$ zu $Z$ wird blockiert, wenn $Y$ beobachtet wird ($X \perp Z \mid Y$).
    \item \textbf{Divergent} $X \leftarrow Y \to Z$: Gemeinsamer Elternknoten; $X$ und $Z$ werden unabhängig, wenn $Y$ beobachtet wird ($X \perp Z \mid Y$).
    \item \textbf{Konvergent} $X \to Y \leftarrow Z$: Gemeinsamer Kindknoten; $X$ und $Z$ sind unabhängig, wenn $Y$ \emph{nicht} beobachtet wird. Beobachten von $Y$ (oder eines Nachfahren) aktiviert Abhängigkeit -- \emph{„explaining away"}.
\end{itemize}}

\definition{\textbf{Naive-Bayes-Klassifikator:} Nimmt bedingte Unabhängigkeit aller Features gegeben die Klasse an: $P(X_1,\ldots,X_p \mid C) = \prod_j P(X_j \mid C)$. Klassifikation via MAP-Schätzung im Log-Raum:
$$\hat{C} = \arg\max_c \left[\log P(C{=}c) + \sum_{j=1}^{p}\log P(X_j \mid C{=}c)\right]$$}

\textbf{Laplace-Smoothing} verhindert Null-Wahrscheinlichkeiten bei ungesehenen Feature-Klasse-Kombinationen:
$$P(X_j{=}x \mid C{=}c) = \frac{\text{count}(X_j{=}x,\, C{=}c) + \alpha}{\text{count}(C{=}c) + \alpha \cdot m}$$
mit Smoothing-Parameter $\alpha$ (typisch $\alpha=1$) und $m$ möglichen Werten von $X_j$.

\textbf{Vorteile:} Schnell, benötigt wenig Daten, liefert Klassifikationssicherheiten, funktioniert überraschend gut bei moderaten Verletzungen der Unabhängigkeitsannahme. \textbf{Nachteile:} Unabhängigkeitsannahme oft unrealistisch; schlechte Kalibrierung bei korrelierten Features.

% ============================================================
\newpage

\subsection{Bayes-Updates: Prior, Likelihood, Evidenz, Posterior}
% ============================================================

\definition{\textbf{Bayesianische Inferenz:} Unbekannte Parameter werden als Zufallsvariablen behandelt. Der Satz von Bayes dient als Update-Regel:
$$p(\theta \mid \mathcal{D}) = \frac{p(\mathcal{D} \mid \theta)\cdot p(\theta)}{p(\mathcal{D})} \qquad \Leftrightarrow \qquad \text{Posterior} = \frac{\text{Likelihood} \times \text{Prior}}{\text{Evidenz}}$$
\begin{itemize}[leftmargin=1.5em]
    \item \textbf{Prior} $p(\theta)$: Vorwissen über $\theta$ vor Beobachtung der Daten; kann uninformativ sein.
    \item \textbf{Likelihood} $p(\mathcal{D}\mid\theta) = \prod_i p(\mathbf{x}_i \mid \theta)$: Wie gut erklärt $\theta$ die Daten?
    \item \textbf{Evidenz} $p(\mathcal{D}) = \int p(\mathcal{D}\mid\theta)\,p(\theta)\,d\theta$: Normalisierungskonstante; relevant für Modellvergleich (Bayes Factor).
    \item \textbf{Posterior} $p(\theta\mid\mathcal{D})$: Aktualisiertes Wissen, balanciert Prior und Likelihood.
\end{itemize}}

\definition{\textbf{Posterior-Punktschätzungen:}
\begin{itemize}[leftmargin=1.5em]
    \item \textbf{MAP} (Maximum A Posteriori): $\hat{\theta}_{\text{MAP}} = \arg\max_\theta\, p(\theta\mid\mathcal{D})$ -- Modus des Posteriors.
    \item \textbf{Posterior Mean}: $\hat{\theta} = \mathbb{E}[\theta\mid\mathcal{D}]$ -- Erwartungswert.
    \item \textbf{Posterior Median}.
\end{itemize}
\textbf{Credible Intervals:} Im Gegensatz zum frequentistischen KI direkt probabilistisch interpretierbar: mit $(1-\alpha)$-Kredibilität liegt $\theta$ tatsächlich im Intervall.}

\definition{\textbf{Sequential Bayesian Update:} Der alte Posterior wird zum neuen Prior:
$$p(\theta \mid \mathcal{D}_1, \mathcal{D}_2) \propto p(\mathcal{D}_2 \mid \theta)\cdot p(\theta \mid \mathcal{D}_1)$$
Die Reihenfolge der Datenpunkte ist unerheblich; Batch- und Online-Update liefern identische Ergebnisse.}

\textbf{Beispiel -- Münzwurf:} Prior $p \sim \text{Beta}(1,1)$ (uninformativ), 5 Würfe, 3 Erfolge. Posterior: $\text{Beta}(4,3)$, Erwartungswert $4/7 \approx 0{,}57$. MLE wäre $3/5 = 0{,}6$; der Bayes-Schätzer ist durch den Prior leicht dazu gezogen.

% ============================================================
\newpage

\subsection{Konjugierte Priors}
% ============================================================

\definition{\textbf{Konjugierter Prior:} Ein Prior $p(\theta)$ heißt konjugiert zur Likelihood $p(\mathcal{D}\mid\theta)$, wenn der Posterior $p(\theta\mid\mathcal{D})$ aus \emph{derselben Verteilungsfamilie} stammt wie der Prior. Vorteil: analytische Update-Regeln ohne numerische Integration.}

\textbf{Wichtige Konjugationen:}
\begin{itemize}[leftmargin=1.5em]
    \item \textbf{Beta--Binomial:} Prior $\text{Beta}(\alpha,\beta)$; nach $k$ Erfolgen in $n$ Versuchen: Posterior $\text{Beta}(\alpha+k,\, \beta+n-k)$. Update-Regel: Erfolge zu $\alpha$, Misserfolge zu $\beta$ addieren.

    \item \textbf{Normal--Normal} (bekannte Varianz $\sigma^2$): Prior $\theta \sim N(\mu_0, \tau_0^2)$; Posterior $N(\mu_n, \tau_n^2)$ mit:
    $$\tau_n^{-2} = \tau_0^{-2} + n\sigma^{-2}, \qquad \mu_n = \tau_n^2\!\left(\tau_0^{-2}\mu_0 + n\sigma^{-2}\bar{x}\right)$$
    Der Posterior ist ein präzisionsgewichteter Durchschnitt von Prior-Mittelwert und Stichprobenmittelwert. Präzisionen addieren sich.

    \item \textbf{Dirichlet--Multinomial:} Prior $\text{Dir}(\alpha_1,\ldots,\alpha_k)$; Posterior $\text{Dir}(\alpha_1+n_1,\ldots,\alpha_k+n_k)$. Anwendung: Naive Bayes, Textklassifikation.
\end{itemize}

\begin{center}
\small
\begin{tabular}{|l|l|l|}
\hline
\textbf{Likelihood} & \textbf{Konjugierter Prior} & \textbf{Posterior-Familie} \\
\hline
Bernoulli/Binomial & Beta & Beta \\
\hline
Multinomial & Dirichlet & Dirichlet \\
\hline
Normal ($\mu$, $\sigma^2$ bekannt) & Normal & Normal \\
\hline
Normal ($\sigma^2$, $\mu$ bekannt) & Inverse-Gamma & Inverse-Gamma \\
\hline
Poisson / Exponential & Gamma & Gamma \\
\hline
\end{tabular}
\end{center}

\definition{\textbf{Normal-Inverse-Gamma (NIW)} -- wenn sowohl $\mu$ als auch $\sigma^2$ unbekannt:
$$p(\mu, \sigma^2) = N\!\left(\mu \,\Big|\, \mu_0,\, \frac{\sigma^2}{\kappa_0}\right) \cdot \text{Inv-Gamma}(\sigma^2 \mid \alpha_0, \beta_0)$$
Update-Regeln nach $n$ Beobachtungen:
$$\kappa_n = \kappa_0 + n, \quad \mu_n = \frac{\kappa_0\mu_0 + n\bar{x}}{\kappa_n}, \quad \alpha_n = \alpha_0 + \frac{n}{2}$$
$$\beta_n = \beta_0 + \frac{1}{2}\sum_i(x_i - \bar{x})^2 + \frac{n\kappa_0}{2\kappa_n}(\bar{x}-\mu_0)^2$$
Der letzte Term in $\beta_n$ bestraft die Abweichung zwischen Stichprobenmittelwert und Prior-Mittelwert.}


% ============================================================

\subsection{Survivorship Bias und kognitive Verzerrungen}
% ============================================================

\definition{\textbf{Survivorship Bias:} Analyse basiert nur auf den „Überlebenden" eines Prozesses; Misserfolge fehlen systematisch, was zu verzerrten Schlussfolgerungen führt:
$$\mathbb{E}[\text{Statistik} \mid D_{\text{obs}}] \neq \mathbb{E}[\text{Statistik} \mid D_{\text{obs}} \cup D_{\text{missing}}]$$}

\textbf{Beispiel -- Flugzeug-Panzerfrage (Wald, 2. WK):} Zurückgekehrte Flugzeuge zeigten Treffer an Tragflächen und Rumpf -- naïve Schlussfolgerung: diese Bereiche panzern. Richtige Schlussfolgerung: Flugzeuge mit Treffern an Cockpit oder Motoren kehrten \emph{nicht} zurück und fehlen in der Analyse. Gerade die Bereiche \emph{ohne} gehäufte Treffer sind die kritischen.

\textbf{Beispiel -- Helm-Paradoxon (1. WK):} Nach Einführung von Stahlhelmen stiegen registrierte Kopfverletzungen im Feldspital. Erklärung: Vorher starben Soldaten mit Kopftreffern sofort und wurden nie erfasst; Helme erhöhten die Überlebensrate, sodass mehr Verletzte ins Spital kamen. Die scheinbare Zunahme misst Überlebende, nicht Gesamtverletzungen.

\definition{\textbf{Availability Bias:} Leicht verfügbare oder medial präsente Ereignisse werden in ihrer Häufigkeit überschätzt (z.\,B. Flugzeugunglücke werden häufiger eingeschätzt als Autounfälle, obwohl umgekehrt).}

\definition{\textbf{Confirmation Bias:} Nur Daten werden erfasst oder bevorzugt, die bestehende Überzeugungen bestätigen -- in ML-Systemen entsteht dadurch eine gefährliche positive Rückkopplungsschleife.

\textbf{Beispiel -- Algorithmischer Bias:} Polizei-Einsatzmodell berechnet Verbrechensdichte aus historischen Daten, weist mehr Polizei in betroffene Viertel zu, erfasst dadurch mehr Kleinkriminalität, erhöht die berechnete Dichte weiter -- self-reinforcing loop. Das System verstärkt die historische Verzerrung statt sie zu korrigieren.}

\definition{\textbf{Gegenmaßnahmen:} Prüfen, welche Daten \emph{fehlen} und warum (Selection Mechanism); Missing Data Analysis und Sensitivitätsanalyse; Blind Randomization; bei gesellschaftskritischen ML-Systemen aktives Fairness-Monitoring, Feedback-Loop-Kontrolle und explizite Fairness Constraints.}

% ============================================================
\newpage

\subsection{Wahrheitsmatrix und abgeleitete Kenngrößen}
% ============================================================

\definition{\textbf{Confusion Matrix:} Tabellarische Darstellung der Klassifikationsleistung:}

\begin{center}
\small
\begin{tabular}{|l|c|c|}
\hline
 & \textbf{Actual Negative} & \textbf{Actual Positive} \\
\hline
\textbf{Predicted Negative} & TN & FN \\
\hline
\textbf{Predicted Positive} & FP & TP \\
\hline
\end{tabular}
\end{center}

FP = Fehlalarm (falsch positiv); FN = übersehen (falsch negativ).

\definition{\textbf{Abgeleitete Metriken:}
\begin{itemize}[leftmargin=1.5em]
    \item \textbf{Accuracy:} $\frac{\text{TP}+\text{TN}}{n}$ -- Gesamtanteil korrekter Vorhersagen; irreführend bei stark unbalancierten Klassen.
    \item \textbf{Sensitivity / Recall / TPR:} $\frac{\text{TP}}{\text{TP}+\text{FN}}$ -- Anteil der echten Positiven, die erkannt wurden. Synonym: Hit Rate.
    \item \textbf{Specificity / TNR:} $\frac{\text{TN}}{\text{TN}+\text{FP}}$ -- Anteil der echten Negativen, die korrekt erkannt wurden.
    \item \textbf{Precision / PPV:} $\frac{\text{TP}}{\text{TP}+\text{FP}}$ -- Von allen vorhergesagten Positiven: wie viele sind korrekt?
    \item \textbf{F1-Score:} $\frac{2\cdot\text{Precision}\cdot\text{Recall}}{\text{Precision}+\text{Recall}} = \frac{2\,\text{TP}}{2\,\text{TP}+\text{FP}+\text{FN}}$ -- Harmonisches Mittel; bestraft extrem niedrige Werte stärker als das arithmetische Mittel. Verallgemeinerung: $F_\beta$ gewichtet Recall stärker ($\beta > 1$) oder Precision stärker ($\beta < 1$).
    \item \textbf{Balanced Accuracy:} $\frac{\text{TPR}+\text{TNR}}{2}$ -- Besser als Accuracy bei unbalancierten Daten.
\end{itemize}}

\textbf{Tradeoffs:} Sensitivity und Specificity sind negativ korreliert -- höherer Threshold $\Rightarrow$ mehr Spezifität, weniger Sensitivität. Wahl hängt vom Kontext ab: Krebsdiagnose bevorzugt hohe Sensitivity (lieber ein falsch-positives als einen Fall übersehen); Spam-Filter bevorzugt hohe Specificity (legitime Emails nicht löschen).

\definition{\textbf{ROC-Kurve und AUC:} Plottet TPR (Y-Achse) gegen FPR $= 1-\text{Specificity}$ (X-Achse) über alle Threshold-Werte. Die Diagonale entspricht einem Zufallsklassifikator; je weiter die Kurve nach oben-links gewölbt, desto besser. \textbf{AUC} (Fläche unter der ROC-Kurve) ist ein threshold-unabhängiges Gütemaß: $0{,}5 =$ Zufall, $1{,}0 =$ perfekt. \textbf{Wahrscheinlichkeitsinterpretation:} AUC $= P(\hat{p}_+ > \hat{p}_-)$ für ein zufälliges positives vs.\ negatives Beispiel.}

\begin{center}
\small
\begin{tabular}{|l|l|l|l|}
\hline
\textbf{Metrik} & \textbf{Formel} & \textbf{Fokus} & \textbf{Geeignet für} \\
\hline
Accuracy & $(\text{TP}+\text{TN})/n$ & Gesamtkorrektheit & Balancierte Daten \\
\hline
Sensitivity/Recall & $\text{TP}/(\text{TP}+\text{FN})$ & Positive finden & Medizin, Sicherheit \\
\hline
Specificity & $\text{TN}/(\text{TN}+\text{FP})$ & Negative erkennen & Spam-Filter \\
\hline
Precision & $\text{TP}/(\text{TP}+\text{FP})$ & Zuverlässigkeit & Kritische Klassifikation \\
\hline
F1-Score & $2PR/(P+R)$ & Balance P \& R & Unbalancierte Daten \\
\hline
AUC & Fläche unter ROC & Ranking-Qualität & Allgemeiner Vergleich \\
\hline
\end{tabular}
\end{center}

% ============================================================
\newpage

\subsection{Support Vector Machines}
% ============================================================

\definition{\textbf{Grundidee:} Finde die Hyperebene $H\colon \mathbf{w}^T\mathbf{x} + b = 0$, die zwei Klassen $y_i \in \{-1,+1\}$ mit maximalem Margin trennt (\emph{Maximum Margin Principle}). Klassifikation: $\hat{y} = \text{sign}(\mathbf{w}^T\mathbf{x}_0 + b)$. Die Lösung hängt nur von den \textbf{Support Vectors} ab -- Trainingspunkten genau auf dem Margin; weiter entfernte Punkte haben keinen Einfluss.}

\definition{\textbf{Hard-Margin SVM} (linear separierbare Daten):
$$\min_{\mathbf{w},b}\;\tfrac{1}{2}\|\mathbf{w}\|^2 \quad \text{s.t.}\quad y_i(\mathbf{w}^T\mathbf{x}_i + b) \geq 1 \;\forall\, i$$
Der Gewichtsvektor ist eine Linearkombination der Trainingspunkte: $\mathbf{w} = \sum_i \alpha_i y_i \mathbf{x}_i$. KKT-Bedingung: $\alpha_i > 0$ genau für Support Vectors.}

\definition{\textbf{Soft-Margin SVM} (Support Vector Classifier, nicht-separierbare Daten): Slack-Variablen $\xi_i \geq 0$ erlauben Constraint-Verletzungen ($\xi_i \leq 1$: in Margin; $\xi_i > 1$: falsch klassifiziert):
$$\min_{\mathbf{w},b,\boldsymbol{\xi}}\;\tfrac{1}{2}\|\mathbf{w}\|^2 + C\sum_i \xi_i \quad \text{s.t.}\quad y_i(\mathbf{w}^T\mathbf{x}_i+b) \geq 1-\xi_i,\; \xi_i \geq 0$$
$C$ steuert den Bias-Varianz-Tradeoff: großes $C$ $\Rightarrow$ enges Margin, Overfitting-Risiko; kleines $C$ $\Rightarrow$ breites Margin, mehr Fehler toleriert. Optimales $C$ via Cross-Validation.}

\definition{\textbf{Kernel Trick:} Statt expliziter Transformation $\phi(\mathbf{x})$ wird die Kernel-Funktion $K(\mathbf{x}_i, \mathbf{x}_j) = \phi(\mathbf{x}_i)^T\phi(\mathbf{x}_j)$ direkt berechnet -- erzeugt nichtlineare Entscheidungsgrenzen ohne explizite Transformation. Gebräuchliche Kernels:
\begin{itemize}[leftmargin=1.5em]
    \item \textbf{Linear:} $K = \mathbf{x}_i^T\mathbf{x}_j$ -- lineare Grenze.
    \item \textbf{Polynomial:} $K = (1+\mathbf{x}_i^T\mathbf{x}_j)^d$ -- polynomiale Grenze.
    \item \textbf{RBF (Gauß):} $K = \exp(-\gamma\|\mathbf{x}_i-\mathbf{x}_j\|^2)$ -- flexibelster, beliebtester Kernel; $\gamma$ steuert Lokalität.
    \item \textbf{Sigmoid:} $K = \tanh(\kappa\mathbf{x}_i^T\mathbf{x}_j+\theta)$.
\end{itemize}}

\textbf{Multi-Klassen-SVM:} \emph{One-vs-One (OvO):} $\binom{K}{2}$ binäre SVMs, Mehrheitsvote. \emph{One-vs-All (OvA):} $K$ binäre SVMs, höchster Konfidenz-Score gewinnt.

\begin{center}
\small
\begin{tabular}{|l|l|l|l|}
\hline
\textbf{Aspekt} & \textbf{SVM} & \textbf{Log.\ Regression} & \textbf{Random Forest} \\
\hline
Hochdimensional & Sehr gut & Gut & Mäßig \\
\hline
Kernel möglich & Ja & Nein & Nein \\
\hline
Rechenaufwand & Hoch ($O(n^3)$) & Niedrig & Mittel \\
\hline
Interpretierbarkeit & Niedrig & Hoch & Mittel \\
\hline
Multi-class nativ & Nein & Ja & Ja \\
\hline
\end{tabular}
\end{center}

% ============================================================
\newpage

\subsection{Validierung, Underfitting und Overfitting}
% ============================================================

\definition{\textbf{Bias-Varianz-Zerlegung:} Der erwartete Testfehler zerfällt in drei Komponenten:
$$\text{MSE} = \underbrace{\text{Bias}^2}_{\text{Underfitting}} + \underbrace{\text{Varianz}}_{\text{Overfitting}} + \underbrace{\sigma^2_\varepsilon}_{\text{irreduzibel}}$$
Bias: systematischer Fehler (Modell zu simpel). Varianz: Empfindlichkeit gegenüber Trainingsdaten (Modell zu komplex). Irreduzibles Rauschen ist nicht eliminierbar.}

\begin{center}
\small
\begin{tabular}{|l|l|l|l|}
\hline
\textbf{Situation} & \textbf{Train-Fehler} & \textbf{Test-Fehler} & \textbf{Maßnahme} \\
\hline
Underfitting & Hoch & Hoch & Komplexeres Modell, mehr Features, weniger Regularisierung \\
\hline
Optimal & Niedrig & Niedrig & -- \\
\hline
Overfitting & Niedrig & Hoch & Regularisierung, mehr Daten, einfacheres Modell \\
\hline
\end{tabular}
\end{center}

\definition{\textbf{Validierungsstrategien:}

\emph{Hold-out:} Aufteilung in Trainings- (60--80\%) und Testset (20--40\%). Einfach, aber hohe Varianz der Schätzung bei kleinen Datensätzen.

\emph{$k$-Fold Cross-Validation:} Daten in $k$ Folds aufteilen; jede Fold dient einmal als Testset, der Rest als Trainingsset:
$$\text{CV-Fehler} = \frac{1}{k}\sum_{i=1}^{k} \text{Error}_i$$
Stabiler als Hold-out, nutzt alle Daten. Typisch: $k=5$ oder $k=10$. LOOCV ($k=n$): sehr rechenintensiv, geringe Bias aber hohe Varianz. \emph{Stratified $k$-Fold} erhält Klassenverteilung in jedem Fold -- wichtig bei unbalancierten Daten.}

\definition{\textbf{Regularisierung:} Reduziert Overfitting durch Bestrafung großer Koeffizienten. L1 (Lasso): $\lambda\sum|w_j|$ -- erzeugt Sparsität (Feature Selection). L2 (Ridge): $\lambda\sum w_j^2$ -- schrumpft Koeffizienten gleichmäßig. Elastic Net: Kombination beider. Hyperparameter $\lambda$ via Cross-Validation wählen: $\lambda=0$ $\Rightarrow$ kein Effekt; großes $\lambda$ $\Rightarrow$ Underfitting.}

\textbf{Early Stopping:} Training abbrechen, sobald Validierungsfehler zu steigen beginnt; Modell mit minimalem Validierungsfehler speichern. Automatische, datengetriebene Modellauswahl ohne zusätzliche Hyperparameter.

\textbf{Learning Curves:} Plot von Train- und Test-Fehler über Modellkomplexität oder Trainingsgröße. Beide hoch und nah beieinander $\Rightarrow$ Underfitting; großer Gap $\Rightarrow$ Overfitting; beide niedrig $\Rightarrow$ optimal.

% ============================================================
\newpage

\subsection{Zeitreihenanalyse}
% ============================================================

\definition{\textbf{Zeitreihe:} Sequenz $(y_t)_{t=1}^T$ in fester zeitlicher Ordnung -- eine Realisierung eines stochastischen Prozesses $(Y_t)$. Beobachtungen sind zeitlich abhängig (nicht i.i.d.), typischerweise äquidistant.}

\definition{\textbf{Stationarität (schwache Form):} Konstanter Mittelwert $E(Y_t)=\mu$, konstante Varianz $\text{Var}(Y_t)=\sigma^2$, und Autokovarianz $\text{Cov}(Y_t,Y_{t+h})$ hängt nur vom Lag $h$ ab, nicht von $t$. Voraussetzung für ARMA-Modelle.}

\definition{\textbf{ACF und PACF:} Empirische Autokorrelation zum Lag $\omega$:
$$r_\omega = \frac{\sum_{t=\omega+1}^{T}(y_t-\bar{y})(y_{t-\omega}-\bar{y})}{\sum_{t=1}^{T}(y_t-\bar{y})^2}$$
Schnell abfallende ACF: Kurzzeitabhängigkeit (stationär); langsam abfallende ACF: nicht-stationär; saisonale Spitzen: periodische Struktur. Die \textbf{PACF} misst Korrelation bei Lag $\omega$ nach Entfernung dazwischenliegender Lags -- hilfreich zur AR-Ordnungsbestimmung.}

\definition{\textbf{Komponentenmodell:} Zerlegung in Trend $\mu_t$, Saisonalität $s_t$ und Restkomponente $\varepsilon_t$. \textbf{Additiv:} $y_t = \mu_t + s_t + \varepsilon_t$ (konstante Saisonamplitude). \textbf{Multiplikativ:} $y_t = \mu_t \cdot s_t \cdot \varepsilon_t$ (proportionale Amplitude; durch Logarithmieren in additives Modell überführbar). Trendschätzung via \textbf{Simple Moving Average}:
$$\hat{\mu}_t = \frac{1}{2q+1}\sum_{i=-q}^{q} y_{t+i}$$
Größeres $q$ $\Rightarrow$ stärkere Glättung; $q$ Randbeobachtungen gehen verloren.}

\definition{\textbf{AR, MA und ARMA-Prozesse:}
\begin{itemize}[leftmargin=1.5em]
    \item \textbf{AR($p$):} $Y_t = \sum_{k=1}^p \phi_k Y_{t-k} + \varepsilon_t$ -- Regression auf $p$ Lags. PACF bricht bei Lag $p$ ab.
    \item \textbf{MA($q$):} $Y_t = \varepsilon_t + \sum_{k=1}^q \theta_k \varepsilon_{t-k}$ -- gewichtete Summe vergangener Schocks. ACF bricht bei Lag $q$ ab.
    \item \textbf{ARMA($p,q$):} Kombination; Voraussetzung: Stationarität.
\end{itemize}}

\definition{\textbf{ARIMA und SARIMA:} \textbf{ARIMA($p,d,q$):} $d$-faches Differenzieren $\Delta^d Y_t$ erzeugt Stationarität aus trendebehafteten Reihen, dann ARMA. \textbf{SARIMA$(p,d,q)\times(P,D,Q)_s$:} Zusätzliche saisonale Komponenten mit Periode $s$.}

\textbf{Modellidentifikation:} PACF bricht bei Lag $p$ ab $\Rightarrow$ AR($p$); ACF bricht bei Lag $q$ ab $\Rightarrow$ MA($q$). Modellwahl via AIC/BIC. Residuen müssen White Noise sein (Ljung-Box-Test, ACF nahe 0).

% ============================================================
\newpage

\subsection{Importance Sampling und MCMC}
% ============================================================

\definition{\textbf{Monte Carlo Integration:} Schätze $E_p[f(x)] = \int f(x)p(x)\,dx$ durch $n$ Stichproben aus $p$: $\hat{E} = \frac{1}{n}\sum_i f(x_i) \xrightarrow{n\to\infty} E_p[f(x)]$ (LLN). Problem: Direktes Sampling aus $p$ oft intraktabel (z.\,B. Bayes-Posterior mit unbekannter Normalisierungskonstante).}

\definition{\textbf{Importance Sampling:} Ziehe aus leicht zu samplender \textbf{Proposal-Verteilung} $q(x)$ und korrigiere mit \textbf{Importance Weights} $w(x) = p(x)/q(x)$:
$$E_p[f(x)] = E_q\!\left[f(x)\frac{p(x)}{q(x)}\right] \approx \sum_i \tilde{w}_i f(x_i), \quad \tilde{w}_i = \frac{w_i}{\sum_j w_j}$$
Anforderung: $q(x)>0$ überall wo $p(x)>0$. Schlechte Proposal $\Rightarrow$ extrem variable Gewichte $\Rightarrow$ wenige Stichproben dominieren (\emph{Degeneracy-Problem}). Die \textbf{Effective Sample Size} $\text{ESS} = (\sum w_i)^2/\sum w_i^2$ misst Effizienz ($\text{ESS}=n$: optimal). In hohen Dimensionen versagt Importance Sampling.}

\definition{\textbf{Markov Chain Monte Carlo (MCMC):} Konstruiere eine Markov Chain mit stationärer Verteilung $p(x)$; nach Burn-in-Phase sind Stichproben approximativ aus $p$. \textbf{Metropolis-Hastings} pro Iteration: (1) schlage $x' \sim q(x'\mid x_t)$ vor; (2) akzeptiere mit
$$\alpha = \min\!\left(1,\frac{p(x')\,q(x_t\mid x')}{p(x_t)\,q(x'\mid x_t)}\right)$$
(3) sonst bleibe bei $x_t$. Detailed Balance garantiert stationäre Verteilung $p$.}

\definition{\textbf{Gibbs Sampling:} Für multivariate $\mathbf{X}=(X_1,\ldots,X_p)$ sample iterativ aus den vollständig bedingten Verteilungen $p(x_j\mid x_{-j})$ -- Spezialfall von Metropolis-Hastings mit $\alpha=1$ (keine Rejections). Vorteil: hohe Akzeptanzrate wenn bedingte Verteilungen analytisch verfügbar.}

\definition{\textbf{Konvergenzdiagnostik:}
\begin{itemize}[leftmargin=1.5em]
    \item \textbf{Trace Plots:} Chain sollte stationär um konstantes Niveau oszillieren (kein Trend, kein Drift).
    \item \textbf{ACF der Chain:} Schnell abfallende ACF $\Rightarrow$ gutes Mixing (unabhängigere Stichproben).
    \item \textbf{Gelman-Rubin $\hat{R}$:} Vergleicht Within- und Between-Chain-Varianz mehrerer paralleler Chains; $\hat{R} < 1{,}1$ deutet auf Konvergenz hin.
    \item \textbf{ESS (mit Autokorrelation):} $\text{ESS} = N/(1 + 2\sum_{h=1}^\infty \rho_h)$ -- berücksichtigt, dass autokorrelierte Stichproben weniger Information tragen.
\end{itemize}}

\begin{center}
\small
\begin{tabular}{|l|l|l|}
\hline
\textbf{Aspekt} & \textbf{Importance Sampling} & \textbf{MCMC} \\
\hline
Unabhängige Samples & Ja & Nein (autokorreliert) \\
\hline
Skalierung (hohe Dim.) & Schlecht & Besser \\
\hline
Normalisierungskonstante & Direkt schätzbar & Nicht direkt \\
\hline
Konvergenzdiagnose & Einfach (ESS) & Aufwendig ($\hat{R}$, ACF) \\
\hline
Ideal wenn & Niedrige Dim., gute Proposal & Hohe Dim., komplexer Posterior \\
\hline
\end{tabular}
\end{center}
\newpage